\documentclass[11pt]{article}
\usepackage[margin=1in]{geometry}
\usepackage{amsmath}
\usepackage{amssymb}
\usepackage{amsthm}
\usepackage{booktabs}
\usepackage{graphicx}
\usepackage{xcolor}
\usepackage{caption}
\usepackage{hyperref}
\usepackage[capitalize,noabbrev,nameinlink]{cleveref}

\newtheorem{proposition}{Proposition}
\newtheorem{remark}{Remark}
\graphicspath{{../figures/}}
\setlength{\emergencystretch}{3em}

% The verified theory section (section_theory.tex) uses natbib-style \citet/\citep.
% natbib is not in the approved package list, so map both onto plain numeric \cite.
\newcommand{\citet}[1]{\cite{#1}}
\newcommand{\citep}[1]{\cite{#1}}

\title{Out-of-Sample Workload Feasibility in CSLAP Layouts:\\
an Empirical Study and a Budgeted-Uncertainty Robust Counterpart}
\author{CSLAP research pipeline}
\date{July 2, 2026}

\begin{document}
\maketitle

\begin{abstract}
We audit the workload-feasibility semantics of the CSLAP evaluation pipeline and re-analyze
every stored layout of the project's instance families (iscf480, iscf10kt, iscfco,
planted-affinity, BERNER). The as-run out-of-sample check compares realized test workload
against a \emph{train-volume-calibrated} ceiling; its bindingness $B$ collapses from
$\approx 1.0$ down to $0.11$--$0.13$ across expanding temporal cuts, so late-cut
``\texttt{wl\_broken}=0'' verdicts are volume slack, not robustness. Under a
volume-normalized ceiling the picture reverses: every temporal fold carries 1--3 of 8
stations above the fair-share ceiling at $k{=}1$ (1--4 over the full closure grid), with
peak load ratios up to 2.52 at $k{=}1$ (2.63 over the full grid; 3.45 in the sparse-train
variants), while the cross-validation folds remain balanced to the exact 10\%-slack edge. A retrospective
Bertsimas--Sim calibration shows the realized overloads sit in the budget-uncertainty sweet
spot: median $\Gamma^{\mathrm{cover}}$ of 8/60 and 19/541 assigned products, deviations
head-concentrated (median Gini 0.63/0.82), and $\Gamma^{\mathrm{bind}}=1$ at most stations.
We derive the exact linear robust counterpart of the workload constraint and implement it
(HiGHS plus a deterministic local search, paired over the temporal cuts). Full-magnitude
protection provably does not fit the 10\%-slack capacity contract (infeasibility
certificates at $\Gamma \ge 4$); ex-ante half-magnitude protection at $\Gamma = 1$ cuts the
out-of-sample peak ratio from 1.77 to 1.59 and the excess workload from 17.4\% to 13.0\%
for a 2.3\% visit price, while uniform ceiling-tightening at the same price buys no peak
reduction at all --- targeted budgeted protection dominates its nominal competitor.
\end{abstract}

\section{Introduction and context}
\label{sec:intro}

The Correlated Storage Location Assignment Problem (CSLAP) assigns products
$p \in P$ to picking stations $s \in S$ (binary $x_{ps}$) so as to minimize the total
number of station visits over a set of orders, subject to a slot-capacity constraint
$\sum_p x_{ps} \le \zeta_s$ and a workload constraint
$\sum_p x_{ps} L_p / V_s \le T_s$ per station. The project's preceding research arc took
the covering-enlargement hypothesis H1 --- replace each order by its pattern-closure and
optimize placement on the enlarged instance to buy out-of-sample robustness \emph{for the
visits objective} --- end-to-end through an exact CPLEX re-validation, and falsified it:
the price of robustness is real and the out-of-sample gain is negative, with confidence
intervals excluding zero.\footnote{Internal reports
\texttt{docs/background/7\_scientific\_report\_iscf\_cplex.md} through
\texttt{docs/background/11\_berner\_decomposition.md} in this repository.}

The same arc, however, surfaced one instability in the data that is \emph{not} an artifact
of the covering mechanism: \emph{demand drift}. The drift-characteristics and order-data
diagnostics\footnote{\texttt{docs/background/8\_drift\_characteristics.md} and
\texttt{docs/background/9\_order\_data\_diagnostics.md}.} showed that per-product line counts move
substantially between a training window and a later operating window, while the placement
pipeline calibrates its workload ceiling once, on the training window, and never revisits
it. This study therefore opens a different channel from H1: not the visits objective, but
\textbf{out-of-sample workload feasibility}. The research question is:

\begin{quote}
\emph{Does a CSLAP layout that satisfies the (hard, train-enforced) workload constraint
remain workload-feasible on later demand, and if not, what deviation magnitudes and
protection budgets would a budgeted-uncertainty robust workload constraint have needed?}
\end{quote}

Contributions:
\begin{itemize}
\item \textbf{Pipeline audit} (\cref{sec:pipeline}): precise semantics of the as-run
feasibility check, and a \emph{bindingness} index $B$ showing the check degenerates into a
volume-calibration test on expanding temporal cuts.
\item \textbf{Volume-normalized reanalysis} (\cref{sec:empirical}): a counterfactual
fair-share ceiling that removes the volume artifact, applied to all stored layouts across
six ISCF-derived families plus the industrial BERNER cell; every recomputed as-run count
reproduces the stored \texttt{wl\_broken} exactly (276 result cells checked across
families), and the train-side constraint held in every one of the $>1{,}000$
station-checks --- so all reported violations are genuinely out-of-sample.
\item \textbf{Retrospective Bertsimas--Sim calibration} (\cref{sec:gamma}): three
deviation-magnitude models (M1 proportional, M2 square-root, M3 train-internal empirical),
fitted coverage constants, and the diagnostics $\Gamma^{\mathrm{cover}}$ /
$\Gamma^{\mathrm{bind}}$ quantifying how large a budget the observed drift asks for and
how small a budget already forces retraining.
\item \textbf{Robust counterpart} (\cref{sec:theory}): the exact linear
Bertsimas--Sim counterpart \cite{bertsimas2004} of the CSLAP workload constraint, in the
row-wise framework of \cite{bental2009robust}, with the Soyster box model
\cite{soyster1973convex} as its $\Gamma \to \zeta_s$ limit.
\item \textbf{Implementation and paired temporal experiment} (\cref{sec:experiment}): the
counterpart implemented on HiGHS with a deterministic local-search incumbent pipeline and
evaluated, paired arm-by-arm, over the five temporal cuts --- yielding infeasibility
certificates for full-magnitude protection under the 10\%-slack contract, a
price--severity frontier for partial insurance, and the head-to-head result that
budgeted protection dominates equal-price uniform ceiling-tightening on out-of-sample
peak severity.
\end{itemize}

\section{The current evaluation pipeline}
\label{sec:pipeline}

\paragraph{Placement (train side).}
Layouts are produced by the Hexaly placement model, which enforces the workload
constraint \emph{hard} on training data:\footnote{\texttt{milp\_synthetic.py}, lines
70--77; the line counts are passed by the harness at lines 427 and 441 of
\texttt{run\_robustness\_iscf.py}.} for every station $s$,
$\sum_{p} (L_p / V_s)\, x_{ps} \le T_s$, where $L_p$ is the product's training-window line
count (\texttt{REAL\_LINES}, distinct orders containing $p$). Our recomputation confirms the
constraint held on the training data for every stored layout of every family
(\texttt{viol\_train}=0 in all 800 iscf10kt station-rows, all 160 gamma-family checks, and
all 1{,}088 secondary station-rows).

\paragraph{Ceiling calibration.}
The time capacity is set once per fold from training volume by the 10\%-slack rule
(\texttt{run\_robustness\_iscf.py}, line 211; same formula in
\texttt{iscf\_adapter.py}, line 194):
\begin{equation}
T_s = \Bigl\lceil\, \sigma \cdot \frac{\sum_{p} L^{\mathrm{tr}}_p}{V_s\,|S|} \,\Bigr\rceil,
\qquad \sigma = 1.10,
\label{eq:slackrule}
\end{equation}
with $V_s$ the station speed ($V_s = 1$, uniform, on all ISCF families;
$\sigma = 1.20$ and heterogeneous per-station ceilings on BERNER).

\paragraph{Evaluation (test side).}
The test instance re-uses the \emph{train} stations file
verbatim.\footnote{\texttt{run\_robustness\_iscf.py}, lines 229--233.} The evaluator
(\texttt{evaluate\_layout\_robust.py}, lines 129--142, driven at
\texttt{run\_robustness\_iscf.py}, lines 490--504) computes the test demand
$\ell^{\mathrm{te}}_p$ of each product as its raw test line count, i.e.\
\texttt{groupby('PRODUCT').size()} over the test orders; the realized
workload $W^{\mathrm{te}}_s = \sum_{p \in \Phi_s} \ell^{\mathrm{te}}_p / V_s$ over the
assigned set $\Phi_s$ of the stored layout (a flat JSON dict
\{\texttt{PROD\_<id>}: \texttt{GARE<n>}\}); and the count
$\texttt{wl\_broken} = |\{s : W^{\mathrm{te}}_s > T_s\}|$ (strict inequality) and
\texttt{cap\_broken} analogously against the slot count. Temporal cuts are 5 expanding-train
splits at train fractions $0.5, 0.6, \dots, 0.9$ by \texttt{ORDER\_ID} rank
(train-early/test-late); the cuts share train mass and are a dispersion proxy, not
independent replicates.

\paragraph{The volume artifact.}
The as-run check thus compares \emph{test-window} workload against a
\emph{train-volume-calibrated} ceiling. Define the \textbf{bindingness} of the check as
\begin{equation}
B \;=\; \frac{\sum_p \ell^{\mathrm{te}}_p}{\sum_{s} V_s\, T_s},
\label{eq:bindingness}
\end{equation}
the ratio of total test volume to total ceiling capacity. With uniform speeds and a flat
$T_s$, a layout balanced exactly to the slack edge (peak load $= \sigma \times$
fair share) crosses $T_s$ if and only if $B > 1/\sigma \approx 0.909$; below that threshold
the check cannot fire on any near-balanced layout, regardless of how unbalanced the demand
mix has become. On the expanding temporal design the test window shrinks as the cut moves
late, so $B$ decays mechanically (from $1.005$ to $0.125$ on iscf480) and the as-run
verdict degenerates into a measurement of window length. Conversely, when the test window
carries \emph{more} volume than the train window (sparse-train and affinity families,
$B > 1$), even a perfectly balanced layout must violate, and the as-run count saturates at
8/8 for reasons that have nothing to do with layout quality.

\section{Empirical study}
\label{sec:empirical}

\subsection{Datasets and stored layouts}
\label{sec:datasets}

\begin{table}[htbp]
\centering
\caption{Instance families audited. All ISCF-derived families have 8 uniform stations
(GARE4--GARE11, $V_s = 1$, $\sigma = 1.10$); BERNER has 24 heterogeneous stations
($\sigma = 1.20$). $k$ is the pattern-closure order of the stored layout
($k = 1$ = nominal).}
\label{tab:datasets}
\footnotesize
\setlength{\tabcolsep}{4pt}
\begin{tabular}{@{}llllll@{}}
\toprule
Family & $|P|$ & $\zeta_s$ & Split design & $k$ grid & Scale (fold 0) \\
\midrule
iscf480 CV        & 480   & 60  & 5 random folds            & 1, 2          & 246.7k / 61.7k orders \\
iscf480 temporal  & 480   & 60  & 5 expanding cuts          & 1, 2          & 154{,}207 / 154{,}208 orders \\
iscf10kt temporal & 4{,}328 & 541 & 5 expanding cuts        & 1, 2, 3, 4, 6 & 5{,}000 / 5{,}000 orders \\
iscf10kt t500/t1000/t2000 & 4{,}328 & 541 & train capped     & 1, 2, 3, 4, 6 & 500--2{,}000 / 1{,}000--5{,}000 \\
iscfco mm/ms      & 1{,}600 & 200 & 5 expanding cuts        & 1, 10         & 7{,}500 / 7{,}500 orders \\
planted affinity  & 800   & 100 & 18 arms (2 modes, 3 $\alpha$, 3 reps) & 1, 10 & 1{,}847 / 3{,}645--3{,}986 \\
BERNER            & 2{,}000 & 9--207 & 1 temporal cut (10/13 wk) & 1--20 & 778{,}788 / 220{,}893 lines \\
\bottomrule
\end{tabular}
\end{table}

\Cref{tab:datasets} lists the families. The affinity scale cell quotes the rep-r0 arms
(unique test orders 3{,}645--3{,}986; over all 18 arms the range is 3{,}645--4{,}000).
The affinity replicate r1 is degenerate as-run
(all 12 layout JSONs are empty, stored rows all-zero), leaving 24 informative arms; BERNER
persisted no layouts, so its audit (\cref{sec:berner}) bounds bindingness through
assignment assumptions rather than recomputing realized loads.

\subsection{Metrics}
\label{sec:metrics}

Let $L^{\mathrm{tr}} = \sum_p L^{\mathrm{tr}}_p$ and
$L^{\mathrm{te}} = \sum_p \ell^{\mathrm{te}}_p$ be total train and test lines, and
$\lambda = L^{\mathrm{tr}} / L^{\mathrm{te}}$ the volume ratio. Consistent with the
notation of \cref{sec:theory}:
\begin{itemize}
\item \textbf{As-run check}: $W^{\mathrm{te}}_s > T_s$; count
$\texttt{viol\_asrun} = \texttt{wl\_broken}$.
\item \textbf{Normalized check}: the volume-normalized load is
$W^{\mathrm n}_s = \lambda\, W^{\mathrm{te}}_s$
(equivalently, the raw load is compared against the counterfactual test-volume ceiling
$T^{\mathrm n}_s = \lceil \sigma L^{\mathrm{te}} / (V_s |S|) \rceil$; the two violation
sets coincide on every audited fold). We report
$\texttt{viol\_norm} = |\{s : W^{\mathrm n}_s > T_s\}|$ and the peak ratio
$\mathrm{ratio}^{\mathrm n} = \max_s W^{\mathrm{te}}_s / T^{\mathrm n}_s$.
\item \textbf{Imbalance}: $\max_s W^{\mathrm{te}}_s \big/ \bigl(L^{\mathrm{te}} / (V_s |S|)\bigr)$,
the peak-to-fair-share ratio; values $\le \sigma = 1.10$ are within the slack rule.
\item \textbf{Bindingness} $B$ as in \cref{eq:bindingness}; volume-binding threshold
$1/\sigma \approx 0.909$.
\item \textbf{Excess fraction} $X = \sum_s (W^{\mathrm n}_s - T_s)_+ \big/ \sum_s W^{\mathrm n}_s$,
the share of (normalized) test workload sitting above the ceiling.
\end{itemize}

\paragraph{Sanity.}
The recomputation reproduces the stored \texttt{wl\_broken} in \emph{all} cells: 20
(split, fold, $k$) cells on iscf480, 100 (variant, fold, $k$) rows on iscf10kt, 136
secondary cells (20 iscfco base + 36 affinity + 80 iscfco demand-shift), and the 20 cells
of the $\Gamma$-calibration re-run; recomputed peak workloads match the stored
\texttt{max\_workload} to $<10^{-6}$ (e.g., iscf480 temporal fold 0, $k=1$:
$86{,}019$ vs $T_s = 45{,}365$). \texttt{cap\_broken}=0 everywhere.

\subsection{Primary instance: iscf480 (CV vs.\ temporal)}
\label{sec:iscf480}

\begin{table}[htbp]
\centering
\caption{iscf480, 5-fold random CV. The as-run anchor \texttt{wl\_broken}=0 everywhere is
reproduced, but $B \approx 0.23$: the test window carries only $\sim$23\% of the volume the
ceiling assumes. Normalized violations exist but are hair-thin (peak ratio $\le 1.016$,
$X \le 0.40\%$). Source: \texttt{results/iscf480/per\_fold.csv}.}
\label{tab:cv480}
\small
\begin{tabular}{@{}lccccccc@{}}
\toprule
Fold & $B$ & as-run $k{=}1$ & as-run $k{=}2$ & norm $k{=}1$ & norm $k{=}2$ & ratio$^{\mathrm n}$ $k{=}1$ & $X$ (\%) $k{=}1$ \\
\midrule
r0f0 & 0.228 & 0 & 0 & \textbf{2} & \textbf{2} & 1.016 & 0.24 \\
r0f1 & 0.227 & 0 & 0 & 0 & 0 & 0.999 & 0.00 \\
r0f2 & 0.227 & 0 & 0 & \textbf{2} & \textbf{1} & 1.005 & 0.13 \\
r0f3 & 0.227 & 0 & 0 & \textbf{3} & \textbf{1} & 1.012 & 0.40 \\
r0f4 & 0.227 & 0 & 0 & \textbf{1} & 0 & 1.000 & 0.00 \\
\bottomrule
\end{tabular}
\end{table}

\begin{table}[htbp]
\centering
\caption{iscf480, 5 expanding temporal cuts. As-run counts merely track the decaying
bindingness $B$; the normalized check reveals 2--4 of 8 stations overloaded in
\emph{every} cut, with peak ratios up to 2.52. Violated-station counts of 8 stations.
Source: \texttt{results/iscf480/per\_fold.csv}.}
\label{tab:temporal480}
\small
\begin{tabular}{@{}lccccccc@{}}
\toprule
Cut & $B$ & \multicolumn{2}{c}{as-run ($k{=}1$ / $k{=}2$)} & \multicolumn{2}{c}{norm ($k{=}1$ / $k{=}2$)} & \multicolumn{2}{c}{ratio$^{\mathrm n}$ ($k{=}1$ / $k{=}2$)} \\
\midrule
r0f0 & 1.005 & \textbf{4} & \textbf{2} & \textbf{3} & \textbf{2} & 1.71 & 1.54 \\
r0f1 & 0.693 & \textbf{2} & \textbf{1} & \textbf{3} & \textbf{4} & 1.91 & 1.71 \\
r0f2 & 0.457 & \textbf{1} & 0 & \textbf{3} & \textbf{3} & 2.25 & 1.98 \\
r0f3 & 0.281 & 0 & 0 & \textbf{2} & \textbf{3} & 2.52 & 2.19 \\
r0f4 & 0.125 & 0 & 0 & \textbf{3} & \textbf{3} & 2.38 & 1.88 \\
\bottomrule
\end{tabular}
\end{table}

\textbf{The CV anchor is a volume artifact.} \Cref{tab:cv480} shows that the headline
``\texttt{wl\_broken}=0 on all CV folds'' says nothing about balance: with
$B \approx 0.227$--$0.228$ against a threshold of $0.909$ (\cref{fig:bindingness}), the
train-calibrated ceiling is
roughly $4\times$ too loose for the CV test volume. The normalized check does surface
violations in 4 of 5 folds at $k{=}1$ --- because Hexaly balances the training workload to
exactly the 10\%-slack edge (imbalance 1.09--1.12 across all (fold, $k$)), any drift tips
the peak station over --- but they are \emph{marginal}: peak ratio at most 1.016, excess
fraction at most $0.40\%$. The honest CV statement is ``balanced to the slack edge, no
material overload.''

\textbf{Temporal drift is severe.} Under the temporal design (\cref{tab:temporal480}),
every cut has 2--3 of 8 stations above the normalized ceiling at $k{=}1$ (2--4 over both
$k$), with peak load
ratios 1.71--2.52 and $X = 11.0$--$30.7\%$ of total test workload above the ceiling;
imbalance reaches 2.77 (cut r0f3, $k{=}1$). Meanwhile the as-run count decays
$4, 2, 1, 0, 0$ in lockstep with $B$ ($1.005 \to 0.125$; \cref{fig:violations}): the late
cuts look ``clean''
as-run exactly where the true overloads are worst. Only cut r0f0 is volume-binding
($B = 1.005$), and there the as-run and normalized checks agree (4 vs 3 stations). The
identity of the peak station changes across cuts (GARE10, GARE4, GARE11, GARE11, GARE6 at
$k{=}1$; \cref{fig:heatmap}; top-3 detail in
\texttt{results/iscf480/top3\_overloaded\_temporal.csv}); every
station holds exactly 60 products in every layout (slot capacity fully binding), so the
overloads are pure demand-mix drift, not slot-count effects. The $k{=}2$ closure layouts
are consistently \emph{less extreme} at the peak (mean peak ratio 1.86 vs 2.15; lower in
all 5 cuts) but remain far infeasible, and do not reduce the violation \emph{count}
(cuts r0f1, r0f3).

\begin{figure}[htbp]
\centering
\includegraphics[width=0.95\textwidth]{fig_violations.pdf}
\caption{iscf480 temporal cuts: as-run vs volume-normalized violated-station counts
(of 8). At $k{=}1$ the as-run count decays $4,2,1,0,0$ while the normalized count stays at
$3,3,3,2,3$; at $k{=}2$: as-run $2,1,0,0,0$ vs normalized $2,4,3,3,3$. The as-run decay
tracks test-window volume, not balance. Source: \texttt{results/iscf480/per\_fold.csv}.}
\label{fig:violations}
\end{figure}

\begin{figure}[htbp]
\centering
\includegraphics[width=0.95\textwidth]{fig_bindingness.pdf}
\caption{Bindingness $B$ of the as-run check across iscf480 folds. Temporal $B$ falls
$1.005 \to 0.125$ as the expanding cut shrinks the test window, while CV $B$ is flat at
$\approx 0.227$; both are judged against the volume-binding threshold
$1/1.10 \approx 0.909$. Only temporal cut 0 is volume-binding; everywhere else the as-run
check is structurally unable to fire on a near-balanced layout.}
\label{fig:bindingness}
\end{figure}

\begin{figure}[htbp]
\centering
\includegraphics[width=0.95\textwidth]{fig_ratio_heatmap.pdf}
\caption{Station-level normalized load ratios $W^{\mathrm{te}}_s / T^{\mathrm n}_s$ for
iscf480 temporal, all 8 stations $\times$ 5 cuts, $k{=}1$ (top) and $k{=}2$ (bottom);
diverging scale centered at 1.0 (ratio $>1$ = overloaded under fair-share volume), all 80
cells annotated. Data range 0.20--2.52 at $k{=}1$; $k{=}2$ peaks at 2.19. The overloaded
station moves across cuts. Source: \texttt{results/iscf480/heatmap\_temporal\_k1.csv} and \texttt{...\_k2.csv}.}
\label{fig:heatmap}
\end{figure}

\subsection{Full-catalogue instance: iscf10kt temporal and sparse-train variants}
\label{sec:iscf10kt}

\begin{table}[htbp]
\centering
\caption{iscf10kt temporal ($k \in \{1,2\}$; full $k$ grid in \cref{tab:app10kt}). Same
two-regime pattern as iscf480. Counts of 8 stations. Source:
\texttt{results/iscf10kt/per\_fold.csv}.}
\label{tab:temporal10kt}
\small
\begin{tabular}{@{}lccccccc@{}}
\toprule
Cut & $B$ & \multicolumn{2}{c}{as-run ($k{=}1$ / $k{=}2$)} & \multicolumn{2}{c}{norm ($k{=}1$ / $k{=}2$)} & ratio$^{\mathrm n}$ $k{=}1$ & $X$ (\%) $k{=}1$ \\
\midrule
0 & 0.786 & \textbf{2} & \textbf{2} & \textbf{2} & \textbf{2} & 1.50 & 10.6 \\
1 & 0.550 & \textbf{1} & 0 & \textbf{2} & \textbf{2} & 1.69 & 15.6 \\
2 & 0.367 & 0 & 0 & \textbf{1} & \textbf{2} & 2.15 & 15.8 \\
3 & 0.224 & 0 & 0 & \textbf{3} & \textbf{1} & 1.93 & 24.6 \\
4 & 0.114 & 0 & 0 & \textbf{3} & \textbf{2} & 2.35 & 25.4 \\
\bottomrule
\end{tabular}
\end{table}

On the full 4{,}328-SKU catalogue over a coverage-guaranteed 10k-order subsample
(\cref{tab:temporal10kt}), the stored as-run anchor ($k{=}1$: $2,1,0,0,0$) is reproduced
exactly (fold 0 peak $2{,}464$ vs $T_s = 1{,}904$, as-run ratio 1.29), and again $B$ falls
$0.79 \to 0.11$, all below the 0.909 threshold --- the fold-4 ceiling is $\approx 8.8\times$
the fair-share test volume. Normalized, \emph{all 100} (variant, fold, $k$) stored layouts
of the iscf10kt families violate on 1--3 of 8 stations; imbalance ranges 1.43--3.80
(always above the 1.10 slack rule) and $X = 6.6$--$37.5\%$.

\begin{table}[htbp]
\centering
\caption{Sparse-train effect (iscf10kt, means over 5 folds $\times$ $k \in \{1,2,3,4,6\}$
unless marked $k{=}1$). The dramatic as-run violation counts at small train samples are a
ceiling artifact ($B \gg 1$: with more test volume than total ceiling capacity even a
perfectly balanced layout must violate); the volume-normalized violation count is flat, and
genuine degradation shows only in the $k{=}1$ imbalance. Source:
\texttt{results/iscf10kt/sparse\_train\_effect.csv}.}
\label{tab:sparse}
\small
\begin{tabular}{@{}lccccccc@{}}
\toprule
Variant & train orders & as-run viol. & $B$ & norm viol. & ratio$^{\mathrm n}$ & $X$ (\%) & imbalance ($k{=}1$) \\
\midrule
t500     & 500    & \textbf{7.56} & 5.06 & \textbf{2.24} & 2.06 & 17.7 & 2.55 \\
t1000    & 1{,}000 & \textbf{6.48} & 2.53 & \textbf{2.16} & 1.94 & 17.0 & 2.34 \\
t2000    & 2{,}000 & \textbf{4.52} & 1.27 & \textbf{2.44} & 1.92 & 18.9 & 2.08 \\
temporal (full) & 7{,}000 & \textbf{0.40} & 0.41 & \textbf{2.16} & 1.86 & 17.1 & 2.12 \\
\bottomrule
\end{tabular}
\end{table}

\Cref{tab:sparse} makes the artifact direction explicit in reverse: at t500 the as-run
check breaks a mean 7.56 of 8 stations ($k{=}1$ by fold: $8,8,8,8,4$), but
$B = 2.07$--$8.18$ there, so those counts measure train-volume calibration, not placement
quality. Normalized, the violation count is essentially flat across train sizes
(2.24 / 2.16 / 2.44 / 2.16), and the genuine sparse-data cost appears only as $k{=}1$
imbalance growing from 2.12 (full) to 2.55 (t500). The single worst layout overall is
t500 fold 4 $k{=}1$: peak ratio$^{\mathrm n}$ 3.45, imbalance 3.80 --- one station carries
47.5\% of test workload.

\subsection{Effect of closure order $k$ on balance}
\label{sec:keffect}

\begin{table}[htbp]
\centering
\caption{$k$-effect on the iscf10kt temporal variant (means over 5 cuts). Closure
enlargement smooths the peak slightly but never approaches feasibility, and the violation
count does not improve monotonically. Source: \texttt{results/iscf10kt/k\_effect.csv}.}
\label{tab:keffect}
\small
\begin{tabular}{@{}lccccc@{}}
\toprule
 & $k{=}1$ & $k{=}2$ & $k{=}3$ & $k{=}4$ & $k{=}6$ \\
\midrule
mean norm violations (of 8) & \textbf{2.2} & \textbf{1.8} & \textbf{1.8} & \textbf{2.4} & \textbf{2.6} \\
mean imbalance              & 2.12 & 2.13 & 2.11 & 2.03 & 1.85 \\
\bottomrule
\end{tabular}
\end{table}

Across all 20 (variant, fold) pairs of the iscf10kt families, \emph{some} $k > 1$ beats
$k{=}1$ on imbalance in 17/20 pairs (mean best-$k$ gain 0.43 imbalance points over the 20
pairs), but the best $k$ is unstable (2, 3, 4 and 6 all occur) and the normalized violation
count never reaches 0 (\cref{tab:keffect}, \cref{fig:keffect}). On iscf480 temporal the
same holds: $k{=}2$ lowers the peak ratio in all 5 cuts (mean 1.86 vs 2.15) without
reducing the violation count. Closures \emph{spread} the overload rather than remove it ---
consistent with the H1 falsification, enlargement is not a workload-feasibility mechanism.

\begin{figure}[htbp]
\centering
\includegraphics[width=0.95\textwidth]{fig_k_effect.pdf}
\caption{$k$-effect panels, iscf10kt temporal variant, $k \in \{1,2,3,4,6\}$: mean
normalized violations stay between 1.8 and 2.6 of 8 stations, and mean imbalance eases only
from 2.12 to 1.85 --- never near the 1.10 slack rule. Source:
\texttt{results/iscf10kt/k\_effect.csv}.}
\label{fig:keffect}
\end{figure}

\subsection{Secondary families and planted shifts}
\label{sec:secondary}

\paragraph{iscfco base (co-occurrence-rich, 1{,}600 SKUs).}
The as-run check reports \texttt{wl\_broken}=0 in 19 of 20 arms (the exception, ms fold 0
$k{=}10$, has 1 station at as-run ratio 1.001), yet 1--3 of 8 stations breach the
normalized ceiling in \emph{every} arm: imbalance 1.12--1.20 (always above 1.10), peak
ratio$^{\mathrm n}$ 1.01--1.09, $X = 0.2$--$2.0\%$. Bindingness collapses $0.858 \to 0.099$
across cuts, so the late-cut zeros are again volume slack. This is the mild edge of the
same phenomenon: material but small overloads on an instance with denser co-occurrence
structure.

\paragraph{Planted demand shift (iscfco, the one mode that moves per-product volumes).}
\Cref{tab:demand} summarizes the only shift family whose stations file carries a
volume-recalibrated ceiling ($B = 0.907$--$0.909 \approx 1/1.10$, i.e., the as-run demand
check is volume-binding by construction, and as-run $\approx$ normalized there). The
nominal $k{=}1$ layout violates on 1--4 of 8 stations at \emph{every} shift intensity
$\rho$, including $\rho = 0$; the violation count is flat in $\rho$ but the severity
scales: mean peak ratio $1.07 \to 1.33$ and $X$ $1.8\% \to 6.9\%$ from $\rho = 0$ to
$\rho = 2$. The $k{=}10$ closure layouts show the same gradient ($X$ $1.0\% \to 5.6\%$) ---
covering enlargement does not repair demand-drift infeasibility.

\begin{table}[htbp]
\centering
\caption{iscfco demand shift, means over 5 cuts (normalized metrics; $k{=}1$ layouts are
identical across the mm/ms variants and deduplicated; $k{=}10$ pooled over mm and ms).
Source: \texttt{results/secondary/per\_fold.csv}.}
\label{tab:demand}
\small
\begin{tabular}{@{}lcccc@{}}
\toprule
$\rho$ & norm viol.\ $k{=}1$ (of 8) & ratio$^{\mathrm n}$ $k{=}1$ & $X$ (\%) $k{=}1$ & $X$ (\%) $k{=}10$ \\
\midrule
0.0 & \textbf{2.6} & 1.07 & 1.8 & 1.0 \\
0.5 & \textbf{2.2} & 1.10 & 2.5 & 2.3 \\
1.0 & \textbf{2.2} & 1.19 & 3.7 & 2.7 \\
2.0 & \textbf{2.6} & 1.33 & 6.9 & 5.6 \\
\bottomrule
\end{tabular}
\end{table}

\paragraph{Planted affinity drift.}
The stored as-run verdicts (7--8 of 8 stations broken in all 24 informative arms, as-run
ratio up to 5.25) are almost entirely volume artifact: test mass is $1.79$--$4.48\times$
the train-calibrated ceiling ($B > 1$). Normalized, the genuine violations are 1--4 of 8
(mean 2.29), imbalance 1.14--1.37, peak ratio$^{\mathrm n}$ up to 1.24 --- and they grow
with planted drift: the reassign arms at $\alpha = 0.5$ / $1.0$ average imbalance
1.318 / 1.285 and $X = 4.8\%$ / $3.6\%$ vs $\approx 1.21$--$1.22$ and $\approx 2\%$ at
$\alpha = 0$. Full per-arm rows in \cref{tab:appaff}.

\paragraph{As-run aggregate over all stored shift evaluations.}
Pooling the 1{,}660 stored rows of the 18 \texttt{per\_fold\_shift\_*.csv} files
(\cref{tab:aggregate}), \texttt{wl\_broken} $> 0$ in 969 rows (58.4\%) and
\texttt{cap\_broken}=0 in all 1{,}660. The split by family confirms the diagnosis: the
sparse-train families are broken in 100\% of rows (ceiling artifact, $B \gg 1$), the
full-train iscf10kt families in $\sim$25--30\% (mean 0.40 of 8), the structured/unstructured
iscfco shifts almost never (mm 0/80, ms 8/80 with max 1, i.e.\ 4/40 in each of the
structured and unstructured modes) --- both shift modes preserve
per-product volumes --- while the demand mode breaks 80/80 rows. As-run counts are
essentially $\rho$-insensitive; only volume calibration and the demand mode matter.

\begin{table}[htbp]
\centering
\caption{As-run aggregate of all stored shift evaluations (per results directory).
``\% rows broken'' = share of rows with \texttt{wl\_broken} $> 0$. Source:
\texttt{results/secondary/shift\_asrun\_aggregate.csv}.}
\label{tab:aggregate}
\small
\begin{tabular}{@{}lcccc@{}}
\toprule
Results dir & rows & \% rows broken & mean \texttt{wl\_broken} & max \\
\midrule
iscf10kt\_results\_mm    & 260 & 30.0  & 0.40 & \textbf{2} \\
iscf10kt\_results\_ms    & 260 & 25.0  & 0.40 & \textbf{2} \\
iscf10kt\_results\_shift & 225 & 28.0  & 0.40 & \textbf{2} \\
iscf10kt\_results\_t500  & 225 & 100.0 & \textbf{7.56} & \textbf{8} \\
iscf10kt\_results\_t1000 & 225 & 100.0 & \textbf{6.48} & \textbf{8} \\
iscf10kt\_results\_t2000 & 225 & 100.0 & \textbf{4.52} & \textbf{8} \\
iscfco\_results\_mm      & 120 & 33.3  & 0.82 & \textbf{4} \\
iscfco\_results\_ms      & 120 & 40.0  & 0.82 & \textbf{4} \\
\bottomrule
\end{tabular}
\end{table}

\subsection{Industrial cell: BERNER audit}
\label{sec:berner}

BERNER is the one heterogeneous family: 24 stations with per-station
$T_s \in \{1, \dots, 158\}$ (calibrated at 120\% of the pinned-assignment train workload,
verified exactly for all 24 stations), speeds 332.6--11{,}885.7, and a single temporal cut
(train quantile 0.7692, 10 of 13 weeks; 778{,}788 raw train lines, 220{,}893 raw test
lines). No layouts were persisted, so realized per-station test loads cannot be
recomputed; the audit instead bounds the check's informativeness four ways, all agreeing:

\begin{itemize}
\item \textbf{Aggregate bindingness}:
$B_{\mathrm{agg}}^{\mathrm{te}} = 220{,}893 / 898{,}472 = 0.246$, far below the
$1/1.20 = 0.833$ threshold (the 3-week test window carries 30.1\% of the distinct-line
\texttt{REAL\_LINES} train volume of $734{,}229$ --- the basis that calibrates $T_s$ ---
and 28.4\% of the 778{,}788 raw train lines). The stored test-demand
\texttt{wl\_broken}=0 for the deployable arms A/B (all 16
decompose rows, every $k$) is therefore \emph{volume-trivial}.
\item \textbf{Assignment bounds}: under the pinned identity assignment the worst station
reaches only ratio 0.286 on test demand (train-frequency-proportional split: 0.251) --- no
train-like layout could plausibly have fired the as-run test check.
\item \textbf{Train side, the contract is binding}:
$B_{\mathrm{agg}}^{\mathrm{tr}} = 734{,}229 / 898{,}472 = 0.817 \approx 1/1.20$ (the
train side of the aggregate bindingness uses the pinned-assignment
\texttt{REAL\_LINES} workload, the quantity $T_s$ was calibrated on; the 778{,}788 raw
train lines would give 0.867), pinned train ratios reach 0.8333
at every large station, and the hard-constrained clean-comparison arms sit at
\texttt{real\_wl\_max\_ratio} 0.998--1.000 --- placement is \emph{at} the ceiling.
\item \textbf{Shape-bindingness}: a fluid equal-workload split of test demand
($W^\ast = 5.674$ per station) would overload station 12 ($T_s = 1$) by $5.67\times$ and
station 24 ($T_s = 2$) by $2.84\times$: equal balancing is incompatible with the
heterogeneous per-station contract, so the constraint binds in shape as well as volume.
\end{itemize}

\begin{table}[htbp]
\centering
\caption{BERNER feasibility semantics by arm group. ``Test'' = TEST raw-line demand vs
real train stations (volume-trivial, $B_{\mathrm{agg}}^{\mathrm{te}} = 0.246$);
``train'' = TRAIN-demand (\texttt{REAL\_LINES}) deployability recheck vs the real
per-station contract. Ranges over the $k$ grid of each source CSV. Source:
\texttt{results/berner/berner\_feasibility\_audit.csv}.}
\label{tab:berner}
\small
\begin{tabular}{@{}llllcc@{}}
\toprule
Arm group & Placement contract & Check & $k$ grid & \texttt{wl\_broken}/24 & max ratio \\
\midrule
A, B (decompose)        & real, hard              & test  & 1--10 & 0 & --- \\
C, D (decompose, recal) & loosened flat            & test  & 2--10 & 0 & --- \\
C$'$, D$'$ (cprime)     & loose recal + $\bar L$ signal & train & 2--10 & \textbf{6--9} & 6.326 \\
B (clean, real regime)  & real, hard              & train & 1--20 & 0 & 1.000 \\
C$''$, D$''$ (fair-inflated) & per-station fair $T^{\mathrm{fair}}_s$ & train & 2--20 & \textbf{4--8} & 1.793 \\
oldD (loose flat cap)   & flat eq.-10 recal        & train & 2--20 & \textbf{7--10} & 2.810 \\
\bottomrule
\end{tabular}
\end{table}

\Cref{tab:berner} separates the two check semantics. The loose-recalibrated C$'$/D$'$ arms
are confirmed non-deployable: 6--9 of 24 stations breach the \emph{real} train-demand
contract, with maximum ratio 5.488 for C$'$ at $k{=}2$ (matching the ``up to $5.5\times$''
figure quoted in the decomposition report) and an overall maximum of 6.326 (D$'$, $k{=}4$).
The hard-constrained A/B arms break 0 stations at every $k$. Note the caveat: BERNER order
CSVs contain duplicate (ORDER, PRODUCT) rows (18{,}367 of 220{,}893 test rows); the as-run
semantics counts raw lines, and this audit follows it. Switching to distinct-order test
demand moves the headline quantities little --- the worst-station pinned test ratio falls
from 0.286 to 0.260 ($-9.2\%$) and the aggregate test volume by 8.3\% --- so no
conclusion changes, though individual small stations move more (the degenerate
$T_s = 1$ station 12 falls from 0.274 to 0.131).

\subsection{Retrospective $\Gamma$ calibration}
\label{sec:gamma}

We instantiate the estimators of \cref{sec:calibration} on the two temporal families
(iscf480, iscf10kt; $k{=}1$ primary, $k{=}2$ secondary), with deviations
$\delta_p = L^{\mathrm n}_p - L^0_p$ of volume-normalized test line counts against the
train nominal ($\lambda$-normalization removes volume drift, so all findings isolate
demand-\emph{mix} drift). The recomputation reproduces every as-run anchor and the
normalized violation sets exactly (no rounding mismatches; 160 train-side checks all
feasible).

\begin{table}[htbp]
\centering
\caption{Fitted deviation constants at coverage $q = 0.90$ (train-line-weighted,
\cref{sec:calibration}), and the as-is coverage of the planner-available M3 (5-block
train-internal std). $c_1$: M1 proportional; $c_2$: M2 square-root. Source:
\texttt{results/gamma/gamma\_fits.csv}.}
\label{tab:fits}
\small
\begin{tabular}{@{}lcccccc@{}}
\toprule
 & \multicolumn{3}{c}{iscf480 temporal} & \multicolumn{3}{c}{iscf10kt temporal} \\
Cut & $c_1(0.9)$ & $c_2(0.9)$ & M3 cov.\ (\%) & $c_1(0.9)$ & $c_2(0.9)$ & M3 cov.\ (\%) \\
\midrule
r0f0 & 0.76 & 46.2  & 5.4 & 1.31 & 5.5  & 79.5 \\
r0f1 & 0.94 & 66.3  & 3.4 & 1.42 & 6.6  & 76.1 \\
r0f2 & 1.18 & 96.5  & 4.2 & 1.63 & 9.7  & 72.0 \\
r0f3 & 1.55 & 146.9 & 1.5 & 2.45 & 13.5 & 58.0 \\
r0f4 & 1.48 & 210.8 & 1.7 & 3.20 & 22.8 & 37.6 \\
\bottomrule
\end{tabular}
\end{table}

\paragraph{Magnitudes (\cref{tab:fits}).}
The fitted constants are large: median $c_1(0.9) = 1.18$ (iscf480) and 1.63 (iscf10kt),
i.e., covering 90\% of train-line-weighted realized deviations requires a deviation
allowance of roughly 118--163\% of a product's nominal line count, and $c_1$ grows with the
cut horizon ($0.76 \to 1.48$ and $1.31 \to 3.20$) --- a single calibration is
horizon-dependent. The train-internal M3, the only genuinely planner-available model,
covers just 1.5--5.4\% of the weighted realized deviations on iscf480 (38--80\% on
iscf10kt): within-train block variability drastically underestimates train$\to$test drift.
These constants are fitted \emph{retrospectively} to realized deviations; they quantify
what would have been needed, not what a planner could estimate ex ante.

\begin{table}[htbp]
\centering
\caption{Budget diagnostics at the $q = 0.90$ fit, $k{=}1$.
$\Gamma^{\mathrm{cover}}$ over the stations violating the normalized check
(iscf480: $n = 14$ station-cells of 60 products each; iscf10kt: $n = 11$ of 541);
$\Gamma^{\mathrm{bind}}$ over all 40 station-cells per family. ``$\infty$'' = uncoverable
even at $\Gamma = |\Phi_s|$ (full box) under the model's magnitudes: for M1/M2 driven by
train-unseen products ($\hat L_p = 0$), for M3 by under-sized block-std magnitudes.
Source: \texttt{results/gamma/gamma\_station.csv}.}
\label{tab:gamma}
\small
\begin{tabular}{@{}llcccccc@{}}
\toprule
 & & \multicolumn{3}{c}{$\Gamma^{\mathrm{cover}}$ (violated stations)} & \multicolumn{3}{c}{$\Gamma^{\mathrm{bind}}$ (all stations)} \\
Family & Model & \#$\infty$ & median & median (finite) & median & \% $=1$ & \% $\le 2$ \\
\midrule
iscf480  & M1 prop.\      & 1  & 4.5      & 4    & 1 & 65.0 & 80.0 \\
iscf480  & M2 sqrt        & 0  & 5        & 5    & 1 & 67.5 & 80.0 \\
iscf480  & M3 empirical   & 10 & $\infty$ & 31   & 3 & 40.0 & 47.5 \\
iscf10kt & M1 prop.\      & 0  & 8        & 8    & 1 & 67.5 & 75.0 \\
iscf10kt & M2 sqrt        & 0  & 28       & 28   & 1 & 57.5 & 70.0 \\
iscf10kt & M3 empirical   & 3  & 135      & 77.5 & 1 & 57.5 & 60.0 \\
\bottomrule
\end{tabular}
\end{table}

\paragraph{Budgets (\cref{tab:gamma}, \cref{fig:gamma}).}
On the stations that actually broke the as-run check, the best-fitting model needs
remarkably small budgets: under M2 on iscf480 the seven as-run-broken stations need
$\Gamma^{\mathrm{cover}} \in \{0,1,2,8,20,23,24\}$, median 8 of 60 assigned products
(13.3\%; M2 is the best mean there, 11.1, with all covers finite, whereas M1 leaves 1/7
uncoverable and M3 4/7); under M1 on iscf10kt they need $\{5, 19, 27\}$, median 19 of 541
(3.5\%, range 0.9--5.0\%). Over all normalized-violated stations the medians remain small
relative to $\zeta_s$ (\cref{tab:gamma}). At the same time $\Gamma^{\mathrm{bind}} = 1$
for 58--68\% of $k{=}1$ stations under M1/M2 (median 1; $\le 2$ for 70--80\%): with these
fitted magnitudes, protecting even the single largest deviation already exceeds the
10\%-slack train ceiling at that majority of stations, so there \emph{any} nonzero budget
forces the optimizer to rebalance at train time (at the remaining stations
$\Gamma^{\mathrm{bind}}$ reaches 109 under M1, so small budgets exert no pressure).
Under M1/M2 the robust constraint would be vacuous nowhere ($\Gamma^{\mathrm{bind}}$
finite in all 80 $k{=}1$ station-cells each); under M3, 12 of 80 station-cells tolerate
any budget ($\Gamma^{\mathrm{bind}} = \infty$).

\paragraph{Concentration (\cref{fig:concentration}).}
Overload is head-concentrated but not degenerate: across the 25 violating (fold, station)
cells at $k{=}1$, the median Gini of positive deviations is 0.63 (iscf480) and 0.82
(iscf10kt); the top-1 product covers 15--50\% (iscf480) and 16--53\% (iscf10kt) of a broken
station's total positive deviation, and the top-10 covers 69--99\% and 61--89\%. A
moderate head of products, not one or two outliers and not diffuse noise, drives the
overloads --- exactly the small-$\Gamma$ regime in which the budget model of
\cref{sec:theory} is designed to operate.

\begin{figure}[htbp]
\centering
\includegraphics[width=0.95\textwidth]{fig_gamma.pdf}
\caption{$\Gamma^{\mathrm{cover}}$ at the $q = 0.90$ fit for $k{=}1$ stations violating
the normalized check, as a fraction of assigned products: iscf480 medians $\approx 7\%$
(M1), $8\%$ (M2), $52\%$ (M3) of $|\Phi_s| = 60$, with 1 (M1) and 10 (M3) infinite covers
annotated; iscf10kt medians $\approx 1.5\%$ (M1), $5\%$ (M2), $14\%$ (M3) of 541, with 3
M3 infinities. Source: \texttt{results/gamma/gamma\_station.csv}.}
\label{fig:gamma}
\end{figure}

\begin{figure}[htbp]
\centering
\includegraphics[width=0.95\textwidth]{fig_concentration.pdf}
\caption{Cumulative share of a broken station's total positive deviation covered by its
top-$m$ deviating products, for the 25 violating (fold, station) cells at $k{=}1$ (both
temporal families). Median cumulative share $\approx 0.29$ at top-1, $\approx 0.61$ at
top-5, $\approx 0.78$ at top-10. Source: \texttt{results/gamma/concentration.csv}.}
\label{fig:concentration}
\end{figure}

% ======================================================================
% Section 4: verified theory section (do not edit its math here)
% ======================================================================
% ======================================================================
% section_theory.tex — Budget-of-uncertainty robust workload constraint
% for the CSLAP placement model. THEORY ONLY (no implementation).
%
% Compilable inside a standard article-class document with
% amsmath, amssymb, amsthm, natbib (authoryear), and cleveref
% (\usepackage[capitalize,noabbrev,nameinlink]{cleveref}, as in the
% project drafts) loaded, and with
%   \newtheorem{proposition}{Proposition}
%   \newtheorem{remark}{Remark}
% defined in the preamble. Citations use \citet/\citep with the keys
% bertsimas2004 (the robust draft's existing key for Bertsimas & Sim 2004;
% do NOT introduce a duplicate bertsimas2004price entry on merge),
% bertsimas2003robust (Bertsimas & Sim 2003, Math. Program.),
% soyster1973convex, bental2009robust.
% The label sec:empirical is wired by the assembler.
% ======================================================================

\section{A budget-of-uncertainty robust counterpart for the workload constraint}
\label{sec:theory}

The empirical companion (\Cref{sec:empirical}) documents a failure mode of
trained CSLAP layouts under demand drift: layouts that satisfy the station
workload constraint on the training window violate it out-of-sample on early
temporal cuts. This section develops the exact robust counterpart of the
workload constraint under budgeted pick-line uncertainty, targeting this
failure mode; how much protection is actually delivered depends on the
calibrated uncertainty set and is assessed in \Cref{sec:calibration}.
Uncertainty is placed on the pick-line coefficients $L_p$ only; the visits
objective is untouched. The construction is the budget-of-uncertainty model
of \citet{bertsimas2004}, in the row-wise robust-optimization framework of
\citet{bental2009robust}, specialized to the CSLAP placement MILP.
Symbols follow the project notation: $P$, $O$, $S$ are the product, order, and
station sets; $P_o$ the support of order $o$; $\zeta_s$, $V_s$, $T_s$ the slot
capacity, processing speed, and time capacity of station $s$.

\subsection{Nominal placement model}
\label{sec:nominal}

% --- removable block: restates the nominal model of the COR manuscript with
% --- its canonical labels (obj, con14--con17) so this section compiles
% --- standalone. Delete on merge into the manuscript and keep the \eqref's.
The nominal placement model assigns every product to exactly one station and
minimizes the total number of station visits over the training orders. The
decision variables are $x_{ps} \in \{0,1\}$ (product $p \in P$ assigned to
station $s \in S$) and the auxiliary visit variables $z_{os} \in [0,1]$ (order
$o \in O$ visits station $s$, completely dependent on $x_{ps}$):
\begin{equation}
\min \sum_{o \in O} \sum_{s \in S} z_{os}
\label{obj}
\end{equation}
subject to
\begin{align}
&\sum_{s \in S} x_{ps} = 1 \quad \forall p \in P, \label{con14}\\
&\sum_{p \in P} x_{ps} \leq \zeta_s \quad \forall s \in S, \label{con15}\\
&x_{ps} \leq z_{os} \quad \forall p \in P_o,\ \forall o \in O,\ \forall s \in S, \label{con16}\\
&\sum_{p \in P} x_{ps} \cdot \frac{L_p}{V_s} \leq T_s \quad \forall s \in S, \label{con17}\\
&x_{ps} \in \{0, 1\}, \quad z_{os} \in [0,1], \quad \forall p \in P,\ o \in O,\ s \in S. \nonumber
\end{align}
Constraint \eqref{con14} assigns each product to exactly one station,
\eqref{con15} is the physical slot capacity, \eqref{con16} links assignments to
visits (a logical OR; no big-$M$ is needed because the objective pushes
$z_{os}$ down to $\max_{p \in P_o} x_{ps}$), and \eqref{con17} bounds the
workload of each station over the planning window.\footnote{The manuscripts typeset the summation in
\eqref{con17} as $\sum_{p \in P_o}$ with no quantifier on $o$; the column-generation
pricing model and the implementation both sum over \emph{all} products assigned
to the station, so we write $\sum_{p \in P}$ here.}

The workload coefficient $L_p$ is a \emph{point estimate}: the training-window
pick-line count of product $p$, computed as the number of distinct training
orders containing $p$ (\texttt{REAL\_LINES}; on instances with duplicate
order--product rows the raw-line and distinct-order conventions differ ---
see the BERNER caveat in \Cref{sec:empirical}). The time capacity is
calibrated on the same window by the
project's $10\%$ operational-slack rule,
\begin{equation}
T_s \;=\; \Bigl\lceil\, 1.10 \cdot \frac{\sum_{p \in P} L_p}{V_s\,|S|} \,\Bigr\rceil,
\label{eq:bs-slack}
\end{equation}
so a layout solving \eqref{obj}--\eqref{con17} is guaranteed workload-feasible
\emph{only for the estimate $L_p$ of the training window}. When the realized
pick-line vector of a later operating window drifts away from the estimate ---
as the temporal experiments of \Cref{sec:empirical} show it does ---
\eqref{con17} offers no protection: a single product whose demand doubles can
push its station past $T_s$ while the model still certifies feasibility.

\subsection{Budget-of-uncertainty model for pick-line workloads}
\label{sec:budget}

Let $L^0_p$ denote the nominal (training-window) estimate of the pick-line
count --- the quantity written $L_p$ in \eqref{con17} --- and let
$\hat L_p \ge 0$ be a deviation magnitude for product
$p$.\footnote{The superscripted symbol $L^0_p$ is chosen deliberately: the
accented $\bar L_p$ is reserved by the covering draft for the
closure-inflated line count (and $\bar T_s$ for its recalibrated capacity),
so no renaming rule is needed when the two studies are merged into one
manuscript.} The uncertain coefficient is modeled as
\begin{equation}
\tilde L_p \in \bigl[L^0_p - \hat L_p,\ L^0_p + \hat L_p\bigr]
\quad \forall p \in P,
\label{eq:bs-uncset}
\end{equation}
with the budget restriction of \citet{bertsimas2004}: in any realization,
\emph{at most $\Gamma$
products deviate from their nominal value simultaneously}, where
$\Gamma \in \{0, 1, \dots, |P|\}$ is the conservatism knob. Since the workload
constraint is a $\le$ row with nonnegative coefficients, only upward
deviations threaten it, so the worst case places every deviating coefficient
at its upper endpoint $L^0_p + \hat L_p$.

Write $\Phi_s = \{p \in P : x_{ps} = 1\}$ for the set of products assigned to
station $s$ (the Hexaly set-variable of the project notation). Requiring
\eqref{con17} to hold for every realization in the budgeted set gives the
robust workload constraint
\begin{equation}
\sum_{p \in \Phi_s} \frac{L^0_p}{V_s}\, x_{ps}
\;+\;
\max_{\substack{J \subseteq \Phi_s \\ |J| \le \Gamma}}\;
\sum_{p \in J} \frac{\hat L_p}{V_s}\, x_{ps}
\;\le\; T_s
\qquad \forall s \in S.
\label{eq:bs-robust}
\end{equation}
On $\Phi_s$ the factors $x_{ps} = 1$ are redundant; they are retained to make
the decision-dependence of the adversary's ground set explicit, and they let
the sums range over all of $P$ without changing value (terms with
$x_{ps} = 0$ vanish). The inner combinatorial maximization admits the
equivalent continuous relaxation
\begin{equation}
\psi_s(x;\Gamma) \;=\; \max_{z} \;
\sum_{p \in P} \frac{\hat L_p}{V_s}\, x_{ps}\, z_p
\quad \text{s.t.} \quad
\sum_{p \in P} z_p \le \Gamma, \qquad
0 \le z_p \le 1 \ \ \forall p \in P.
\label{eq:bs-inner}
\end{equation}
For fixed $x$, \eqref{eq:bs-inner} is a linear program in $z$ alone: the $x$
enter only through the constant objective coefficients
$d_{ps} := (\hat L_p / V_s)\, x_{ps} \ge 0$. Its constraint matrix --- one
cardinality row plus the identity box --- is an interval matrix, hence totally
unimodular; equivalently, by the greedy argument, since all $d_{ps} \ge 0$ an
optimal $z$ sets $z_p = 1$ on the $\Gamma$ largest coefficients (on all
positive ones if fewer than $\Gamma$) and $z_p = 0$ elsewhere. For integer
$\Gamma$ integrality of the inner adversarial problem is therefore
\emph{free}, and $\psi_s(x;\Gamma)$ equals the set-maximum in
\eqref{eq:bs-robust}.\footnote{The adversary variable $z_p$ (single subscript)
is unrelated to the visit variable $z_{os}$ and to the cover-pricing $z_p$ of
the covering draft; it exists only inside \eqref{eq:bs-inner} and is dualized
away in \Cref{prop:bs-counterpart}.}

\begin{proposition}[Exact linear counterpart]
\label{prop:bs-counterpart}
Fix a station $s \in S$, an integer budget $\Gamma \in \mathbb Z_{\ge 0}$,
and a binary assignment vector
$x_{\cdot s} \in \{0,1\}^{P}$.\footnote{Integrality of $\Gamma$ is what makes
the LP value $\psi_s(x;\Gamma)$ of \eqref{eq:bs-inner} coincide with the
set-maximum in \eqref{eq:bs-robust}; for fractional $\Gamma$ the
($\Rightarrow$) direction fails: one assigned product with
$\hat L_p / V_s = 1$, zero nominal slack, and $\Gamma = 1/2$ satisfies
\eqref{eq:bs-robust} (only $J = \emptyset$ is admissible), yet
$\psi_s(x;1/2) = 1/2 > 0$, so no multipliers exist. Real-valued budgets can
be recovered by adding the Bertsimas--Sim fractional deviation term
($\lfloor \Gamma \rfloor$ full deviations plus one scaled by
$\Gamma - \lfloor \Gamma \rfloor$) to \eqref{eq:bs-robust}; integer budgets
suffice here, consistent with $\Gamma \in \{0, 1, \dots, |P|\}$ above.}
The robust workload constraint
\eqref{eq:bs-robust} holds if and only if there exist multipliers
$\theta_s \ge 0$ and $\mu_{ps} \ge 0$, $p \in P$, such that
\begin{align}
&\sum_{p \in P} \frac{L^0_p}{V_s}\, x_{ps}
\;+\; \Gamma\, \theta_s
\;+\; \sum_{p \in P} \mu_{ps}
\;\le\; T_s, \label{eq:bs-cpt-load}\\[2pt]
&\theta_s + \mu_{ps} \;\ge\; \frac{\hat L_p}{V_s}\, x_{ps}
\qquad \forall p \in P. \label{eq:bs-cpt-dualfeas}
\end{align}
Since $\hat L_p / V_s$ and $L^0_p / V_s$ are data constants, each right-hand
side is a constant multiple of the decision variable $x_{ps}$, so
\eqref{eq:bs-cpt-load} and \eqref{eq:bs-cpt-dualfeas} are linear in
$(x, \theta, \mu)$; no integrality of $x$ is needed for linearity.
\end{proposition}

\begin{proof}
By the discussion following \eqref{eq:bs-inner} --- the step in which
integrality of $\Gamma$ is used --- constraint
\eqref{eq:bs-robust} is equivalent to
\begin{equation}
\psi_s(x;\Gamma) \;\le\; T_s - \sum_{p \in P} \frac{L^0_p}{V_s}\, x_{ps}
\;=:\; \rho_s(x),
\label{eq:bs-slackform}
\end{equation}
where $\rho_s(x)$ is the nominal slack of station $s$. The inner LP
\eqref{eq:bs-inner}, with coefficients
$d_{ps} = (\hat L_p / V_s)\, x_{ps} \ge 0$, is \emph{feasible} ($z = 0$) and
\emph{bounded above} ($\sum_p d_{ps} z_p \le \sum_p d_{ps} < \infty$ since
$z_p \le 1$), so it attains a finite optimum $\psi_s(x;\Gamma)$. Associate
the dual variable $\theta_s \ge 0$ with the budget row
$\sum_p z_p \le \Gamma$ and $\mu_{ps} \ge 0$ with each upper bound
$z_p \le 1$. The LP dual of \eqref{eq:bs-inner} is
\begin{equation}
\min_{\theta_s,\, \mu_{\cdot s}} \;\; \Gamma\,\theta_s + \sum_{p \in P} \mu_{ps}
\quad \text{s.t.} \quad
\theta_s + \mu_{ps} \ge d_{ps} \ \ \forall p \in P, \qquad
\theta_s \ge 0,\ \mu_{ps} \ge 0.
\label{eq:bs-dual}
\end{equation}
Because the primal is feasible and bounded, LP strong duality applies: the
dual \eqref{eq:bs-dual} is solvable and its optimal value equals
$\psi_s(x;\Gamma)$. We prove the two implications separately.

($\Leftarrow$) Suppose $(\theta_s, \mu_{\cdot s}) \ge 0$ satisfies
\eqref{eq:bs-cpt-load}--\eqref{eq:bs-cpt-dualfeas}. Constraint
\eqref{eq:bs-cpt-dualfeas} says the point is feasible for the dual
\eqref{eq:bs-dual}, so by \emph{weak} duality
$\psi_s(x;\Gamma) \le \Gamma\,\theta_s + \sum_p \mu_{ps}$. Rearranging
\eqref{eq:bs-cpt-load} gives
$\Gamma\,\theta_s + \sum_p \mu_{ps} \le \rho_s(x)$. Chaining the two
inequalities yields $\psi_s(x;\Gamma) \le \rho_s(x)$, which is
\eqref{eq:bs-slackform}, i.e., the robust constraint \eqref{eq:bs-robust}
holds.

($\Rightarrow$) Conversely, suppose \eqref{eq:bs-robust} holds, i.e.,
$\psi_s(x;\Gamma) \le \rho_s(x)$. By strong duality the dual
\eqref{eq:bs-dual} attains its optimum at some
$(\theta_s^\ast, \mu_{\cdot s}^\ast) \ge 0$ with
\[
\Gamma\,\theta_s^\ast + \sum_{p \in P} \mu_{ps}^\ast
\;=\; \psi_s(x;\Gamma) \;\le\; \rho_s(x)
\;=\; T_s - \sum_{p \in P} \frac{L^0_p}{V_s}\, x_{ps}.
\]
The optimal dual point is dual-feasible, so it satisfies
\eqref{eq:bs-cpt-dualfeas}; the displayed inequality is exactly
\eqref{eq:bs-cpt-load}. Hence the required multipliers exist.

Finally, linearity: $\hat L_p/V_s$ and $L^0_p/V_s$ are constants, so
\eqref{eq:bs-cpt-load} is linear in $(x,\theta,\mu)$ and the right-hand side
of \eqref{eq:bs-cpt-dualfeas} is a constant multiple of the decision variable
$x_{ps}$ --- no product of decision variables appears anywhere, whether or
not $x$ is integral.
\end{proof}

\subsection{The robust placement MILP}
\label{sec:robust-milp}

Applying \Cref{prop:bs-counterpart} station by station --- with a per-station
integer budget $\Gamma_s \in \mathbb Z_{\ge 0}$, the uniform choice
$\Gamma_s = \Gamma$ being the default --- and substituting the counterpart
for \eqref{con17} yields the full
robust placement MILP:
\begin{subequations}
\label{eq:bs-milp}
\begin{align}
&\min_{x,\,z,\,\theta,\,\mu} \; \sum_{o \in O} \sum_{s \in S} z_{os}
  \label{eq:bs-milp-obj}\\
\text{s.t.}\;\;
&\sum_{s \in S} x_{ps} = 1
  \quad \forall p \in P, \label{eq:bs-milp-assign}\\
&\sum_{p \in P} x_{ps} \le \zeta_s
  \quad \forall s \in S, \label{eq:bs-milp-slot}\\
&x_{ps} \le z_{os}
  \quad \forall p \in P_o,\ \forall o \in O,\ \forall s \in S, \label{eq:bs-milp-link}\\
&\sum_{p \in P} \frac{L^0_p}{V_s}\, x_{ps}
  + \Gamma_s\, \theta_s
  + \sum_{p \in P} \mu_{ps} \le T_s
  \quad \forall s \in S, \label{eq:bs-milp-wl}\\
&\theta_s + \mu_{ps} \ge \frac{\hat L_p}{V_s}\, x_{ps}
  \quad \forall p \in P,\ \forall s \in S, \label{eq:bs-milp-dualfeas}\\
&x_{ps} \in \{0,1\},\quad z_{os} \in [0,1],\quad
  \theta_s \ge 0,\quad \mu_{ps} \ge 0
  \quad \forall p,\ o,\ s. \label{eq:bs-milp-dom}
\end{align}
\end{subequations}
The objective \eqref{eq:bs-milp-obj} and the combinatorial core
\eqref{eq:bs-milp-assign}--\eqref{eq:bs-milp-link} are identical to the
nominal model \eqref{obj}, \eqref{con14}--\eqref{con16}; only the workload
block changes. The size overhead over the nominal MILP is exactly $+|S|$
continuous variables $\theta_s$, $+|P|\,|S|$ continuous variables $\mu_{ps}$,
and $+|P|\,|S|$ dual-feasibility constraints \eqref{eq:bs-milp-dualfeas};
the workload rows stay at $|S|$, the binaries and the objective are
unchanged. Robustness is bought with one continuous LP block, not with
additional integer structure.

\subsection{Properties}
\label{sec:properties}

\begin{proposition}[Budget limits and monotonicity]
\label{prop:bs-limits}
For an integer budget $\Gamma \in \mathbb Z_{\ge 0}$, let
\[
\mathcal A(\Gamma) \;=\; \bigl\{ x \in \{0,1\}^{P \times S} \;:\;
\exists\, (z, \theta, \mu) \text{ such that } (x, z, \theta, \mu)
\text{ is feasible for \eqref{eq:bs-milp} with } \Gamma_s = \Gamma \bigr\}
\]
denote the projection onto the assignment variables of the feasible set of
\eqref{eq:bs-milp} with uniform budget $\Gamma$. Then:
\begin{enumerate}
\item[(i)] $\Gamma = 0$ recovers the nominal model: $\mathcal A(0)$ equals
the projection onto $x$ of the feasible set of \eqref{con14}--\eqref{con17}
with $L_p = L^0_p$.
\item[(ii)] If $\Gamma \ge \zeta_s$ for every $s$, then \eqref{eq:bs-robust}
reduces to the box (Soyster) counterpart
$\sum_{p \in P} \frac{L^0_p + \hat L_p}{V_s}\, x_{ps} \le T_s$
\citep{soyster1973convex}.
\item[(iii)] $\Gamma \mapsto \mathcal A(\Gamma)$ is nonincreasing:
$\Gamma' \ge \Gamma$ implies $\mathcal A(\Gamma') \subseteq \mathcal A(\Gamma)$,
and the optimal visit count of \eqref{eq:bs-milp} is nondecreasing in $\Gamma$.
\end{enumerate}
\end{proposition}

\begin{proof}
(i) Both projections are cut out by the common core
\eqref{con14}--\eqref{con16}; the visit variables never restrict $x$ further
(setting $z_{os} = 1$ is always feasible), so it suffices to compare the
workload blocks. With $\Gamma = 0$ the inner maximum in \eqref{eq:bs-robust}
is over $J = \emptyset$ only and equals $0$, so \eqref{eq:bs-robust} is
\eqref{con17}. In the counterpart, given $x$ satisfying \eqref{con17}, choose
$\theta_s = \max_p \hat L_p / V_s$ and $\mu_{ps} = 0$:
\eqref{eq:bs-cpt-dualfeas} holds and $\Gamma_s \theta_s$ vanishes, so
\eqref{eq:bs-cpt-load} is exactly \eqref{con17}; conversely
\eqref{eq:bs-cpt-load} with $\Gamma_s = 0$ and $\mu \ge 0$ implies
\eqref{con17}.

(ii) By the slot capacity \eqref{con15}, $|\Phi_s| = \sum_p x_{ps} \le
\zeta_s \le \Gamma$, so the restriction $|J| \le \Gamma$ is vacuous and the
adversary deviates \emph{every} assigned product upward: the inner maximum is
$\sum_{p \in \Phi_s} \hat L_p / V_s$, the worst case of the full interval box
\eqref{eq:bs-uncset}. Any $\Gamma \ge \max_s |\Phi_s|$ suffices; since
$|\Phi_s|$ is decision-dependent, $\Gamma \ge \zeta_s$ is the natural
data-level sufficient condition.

(iii) $\psi_s(x;\Gamma)$ is nondecreasing in $\Gamma$ (the feasible set of
\eqref{eq:bs-inner} grows with $\Gamma$), so the feasible sets
$\mathcal A(\Gamma)$ are nested and minimizing the same objective over a
smaller set cannot decrease the optimum.
\end{proof}

Thus $\Gamma$ interpolates monotonically between the nominal model and full
box conservatism, exactly the Bertsimas--Sim design.

\begin{remark}[Decision-dependent adversary support]
\label{rem:bs-endogenous}
The adversary's ground set $\Phi_s$ in \eqref{eq:bs-robust} is itself a
decision, since $\Phi_s = \{p : x_{ps} = 1\}$. Endogenous uncertainty sets
are in general a hard extension of robust optimization, but no extra
machinery is needed here: because the factor $x_{ps}$ is kept inside the
inner objective of \eqref{eq:bs-inner}, deviating a non-assigned product
contributes zero, so the adversary restricted to $\Phi_s$ and the adversary
ranging over all of $P$ attain the same value. The dualization of
\Cref{prop:bs-counterpart} then applies verbatim. Linearity of the
counterpart is automatic because $\hat L_p / V_s$ are data constants (no
bilinear terms, no big-$M$); what binariness of $x_{ps}$ buys is this
ground-set equivalence, not linearity. (Under merely nonnegative $x_{ps}$
the equivalence holds with respect to the support $\{p : x_{ps} > 0\}$
rather than $\Phi_s$: a fractional $0 < x_{ps} < 1$ lies outside $\Phi_s$
yet contributes $(\hat L_p / V_s)\, x_{ps} > 0$ to \eqref{eq:bs-inner}.)
\end{remark}

\begin{remark}[Probabilistic guarantee]
\label{rem:bs-prob}
Suppose the realized deviations of the assigned products are independent and
symmetrically distributed on the intervals \eqref{eq:bs-uncset}. Then, for
any assignment $x$ satisfying the protected row \eqref{eq:bs-robust} with
budget $\Gamma$ --- in particular, any $x$ feasible for \eqref{eq:bs-milp}
--- \citet{bertsimas2004} bound the violation probability of that row by
\begin{equation}
\Pr\Bigl[\textstyle\sum_{p \in \Phi_s} \tilde L_p / V_s > T_s\Bigr]
\;\le\; B(n, \Gamma) \;\le\; \exp\!\bigl(-\Gamma^2 / (2n)\bigr),
\qquad n = |\Phi_s| \le \zeta_s,
\label{eq:bs-probbound}
\end{equation}
where $B(n, \Gamma)$ is the exact binomial bound of Theorem~2 of
\citet{bertsimas2004}, well approximated by
$1 - \Phi_{\mathcal N}\bigl((\Gamma - 1)/\sqrt{n}\bigr)$ with
$\Phi_{\mathcal N}$ the standard normal c.d.f.\ (not the assigned set
$\Phi_s$); the closed-form exponential expression is the weaker
Hoeffding-type bound of \citet{bertsimas2003robust}, derived for homogeneous
deviation magnitudes across the row. A budget of order $\sqrt{n}$ already
yields a strong guarantee at a small fraction of the box protection level.
One caveat: daily pick-line counts are nonnegative and right-skewed, so the
symmetric-deviation assumption behind \eqref{eq:bs-probbound} is an
approximation here; the bound is indicative, and in practice it is replaced
by the data-driven procedure of \Cref{sec:calibration}: deviation magnitudes
fitted on train-internal blocks, a budget fixed ex ante on a grid, and
retrospective diagnostics assessing adequacy.
\end{remark}

\begin{remark}[Constraint-side uncertainty only; relation to H1]
\label{rem:bs-objective}
Uncertainty enters model \eqref{eq:bs-milp} through the \emph{constraints}
only; the visits objective \eqref{eq:bs-milp-obj} is identical to the nominal
one. This is deliberately complementary to the falsified covering hypothesis
H1, which perturbed the \emph{objective} by enlarging every order to its
pattern-closure and paid a certified out-of-sample price ($G(\kappa) < 0$ at
$\kappa = 2$ on the clean instance). Here nothing is insured about the visit
count: the insured quantity is \emph{workload-feasibility risk} --- the
out-of-sample violations $W^{\mathrm n}_s > T_s$ documented in
\Cref{sec:empirical} ---
which the nominal model cannot see by construction. The two mechanisms are
orthogonal; the present study evaluates the workload channel alone.
\end{remark}

\begin{remark}[Volume basis and reporting convention]
\label{rem:bs-basis}
The counterpart is meaningful only if $L^0_p$, $\hat L_p$, and $T_s$ are
expressed on the \emph{same volume basis}. In the project's adapters, $T_s$
derives from the training-window totals by the $10\%$ slack rule
\eqref{eq:bs-slack} ($\kappa$-invariant on the clean instance); accordingly,
$L^0_p$ must be the training-window line count and $\hat L_p$ a deviation
of volume-normalized counts on that basis (\Cref{sec:calibration}). Mixing
bases --- e.g., a test-window $\hat L_p$ against a train-calibrated $T_s$ ---
would re-introduce the recalibration confound the clean instance was built to
remove. Since the protection level depends on $n = |\Phi_s| \le \zeta_s$, we
recommend reporting budgets as the dimensionless fraction
$\Gamma_s / \zeta_s$, comparable across stations and instances.
\end{remark}

\subsection{Data-driven calibration of $(\hat L_p, \Gamma)$}
\label{sec:calibration}

The counterpart of \Cref{sec:robust-milp} is fully specified once the
deviation magnitudes $\hat L_p$ and the budget $\Gamma$ are chosen. We define
the estimators the empirical section instantiates; no implementation is
discussed.

\paragraph{Deviation models.}
Let $\delta_p^{(b)} = L_p^{(b)} - L^0_p$ denote the realized deviation of
the volume-normalized line count $L_p^{(b)}$ of product $p$ on a
train-internal temporal block $b \in \{1,\dots,B\}$ (contiguous segments of
the training window by order rank; counts are rescaled so each block carries
the training-window total volume, keeping everything on the basis of
\Cref{rem:bs-basis}). The three deviation models are
\begin{equation}
\text{M1:}\ \ \hat L_p = c_1\, L^0_p,
\qquad
\text{M2:}\ \ \hat L_p = c_2\, \sqrt{L^0_p},
\qquad
\text{M3:}\ \ \hat L_p = \operatorname{std}_{b=1,\dots,B}\bigl(L_p^{(b)}\bigr).
\label{eq:bs-devmodels}
\end{equation}
M1 assumes deviations proportional to demand (multiplicative drift), M2
assumes the $\sqrt{L^0_p}$ scaling of count noise (Poisson-like), and M3
is the nonparametric per-product standard deviation across the training
blocks. M3 needs no fitted constant; M1 and M2 are calibrated as follows.

\paragraph{Fitting $c$.}
For a coverage target $q \in (0,1)$ (e.g.\ $q = 0.9$), the constant is the
smallest value whose implied intervals cover a $L^0$-weighted fraction $q$
of the realized deviations on the train-internal blocks:
\begin{equation}
c(q) \;=\; \min \Bigl\{ c \ge 0 \;:\;
\sum_{p,\,b} L^0_p\, \mathbf 1\bigl[\,|\delta_p^{(b)}| \le c\, g(L^0_p)\bigr]
\;\ge\; q \sum_{p,\,b} L^0_p \Bigr\},
\label{eq:bs-cfit}
\end{equation}
with $g(L^0_p) = L^0_p$ for M1 and $g(L^0_p) = \sqrt{L^0_p}$ for
M2; equivalently, $c(q)$ is the $L^0$-weighted $q$-quantile of the
normalized deviations $|\delta_p^{(b)}| / g(L^0_p)$. Weighting by
$L^0_p$ prioritizes high-volume products, whose deviations dominate
station workloads; restricting to train-internal blocks keeps the calibration
free of test leakage.

\paragraph{Budget diagnostics.}
Fix a trained layout $a$ (indicator $x^a$), a station $s$ with assigned set
$\Phi_s$, nominal load $\bar W_s = \sum_{p \in \Phi_s} L^0_p / V_s$, and
realized normalized load
$W^{\mathrm n}_s = \sum_{p \in \Phi_s} L^{\mathrm n}_p / V_s$, where
$L^{\mathrm n}_p$ is the volume-normalized realized line count of the
evaluation window. By \eqref{eq:bs-inner},
$\psi_s(x^a;\Gamma) = \frac{1}{V_s} \sum_{p \in J_\Gamma(s)} \hat L_p$ with
$J_\Gamma(s) \subseteq \Phi_s$ the $\min(\Gamma, |\Phi_s|)$ largest
deviations at station $s$. Two \emph{retrospective} diagnostics summarize
how large a budget the data ask for and how large a budget the training
instance tolerates:
\begin{align}
\Gamma^{\mathrm{cover}}(s) &\;=\;
\min \Bigl\{ \Gamma \in \mathbb Z_{\ge 0} \;:\;
\bar W_s + \psi_s(x^a;\Gamma) \;\ge\; W^{\mathrm n}_s \Bigr\},
\label{eq:bs-gcover}\\
\Gamma^{\mathrm{bind}}(s) &\;=\;
\min \Bigl\{ \Gamma \in \mathbb Z_{\ge 0} \;:\;
\bar W_s + \psi_s(x^a;\Gamma) \;>\; T_s \Bigr\}.
\label{eq:bs-gbind}
\end{align}
$\Gamma^{\mathrm{cover}}(s)$ is the smallest budget whose protection level
reaches the realized load of station $s$: the budget that would have been
needed for the robust constraint to ``see'' the observed drift (it is $0$
when $\bar W_s \ge W^{\mathrm n}_s$, and undefined --- reported as $\infty$
--- when even $\Gamma = |\Phi_s|$ under-covers, i.e., when the deviation
model itself is too small). $\Gamma^{\mathrm{bind}}(s)$ is the smallest
budget at which the robust left-hand side exceeds $T_s$ on the training data,
i.e., at which the counterpart starts to cut off the trained layout; budgets
below it exert no retraining pressure at station $s$, budgets at or above it
force the optimizer to rebalance. Symmetrically to
$\Gamma^{\mathrm{cover}}$, the set in \eqref{eq:bs-gbind} may be empty: when
even the full box load keeps the station within capacity,
$\bar W_s + \sum_{p \in \Phi_s} \hat L_p / V_s \le T_s$, no finite budget
makes the constraint bind ($\psi_s$ saturates at $\Gamma = |\Phi_s|$), and
we set $\Gamma^{\mathrm{bind}}(s) = \infty$ --- such stations tolerate any
budget.

\paragraph{Use in the empirical study.}
In a future implementation phase the operative budget would be fixed
\emph{ex ante}: robust layouts would be trained on a predeclared grid of
budgets (reported as fractions of $\zeta_s$ per \Cref{rem:bs-basis}), with
the deviation magnitudes fitted on train-internal blocks only via
\eqref{eq:bs-cfit}, so that no evaluation-window quantity would enter the
choice of $\Gamma$ for any scored test cut. No robust layout is trained in
this study: \Cref{sec:empirical} computes only the \emph{retrospective}
diagnostics \eqref{eq:bs-gcover}--\eqref{eq:bs-gbind} on the existing
nominal layouts --- for each temporal cut and deviation model in
\eqref{eq:bs-devmodels}, the distributions of $\Gamma^{\mathrm{cover}}(s)$
over the stations that actually violate $W^{\mathrm n}_s > T_s$
out-of-sample and of $\Gamma^{\mathrm{bind}}(s)$ over \emph{all} stations
(healthy and violated alike) --- answering ``what budget \emph{would have
been} needed'' and ``what budget the training instance tolerates''. The
empirical instantiation also departs from \eqref{eq:bs-cfit} on the
magnitude side: the reported constants $c_1, c_2$ (\Cref{sec:gamma}) apply
the same $L^0$-weighted quantile formula \emph{retrospectively} to the
realized evaluation-window deviations
$\delta_p = L^{\mathrm n}_p - L^0_p$ --- they are diagnostic, not
planner-available --- and only M3 uses train-internal blocks as
\eqref{eq:bs-cfit} prescribes. A budget $\Gamma$ would be judged
\emph{adequate in retrospect} when it dominates the observed
$\Gamma^{\mathrm{cover}}$ quantiles on drift-hit stations --- so the
protected constraint would have flagged the realized violations, and the
robust model is designed to rebalance them away at train time --- while
staying below $\Gamma^{\mathrm{bind}}$ on most healthy stations, so
protection is bought where drift occurs rather than by uniform slack. Once
an ex-ante grid exists, these curves would audit it after the fact rather
than select the budget: they replace the a priori appeal to the
symmetric-noise bound \eqref{eq:bs-probbound} as the lens for \emph{judging}
$\Gamma$, and they must not feed back into its choice on any window used
for scoring.


% ======================================================================
% Section 5: computational study (implementation + paired experiment)
% ======================================================================
\section{Computational study: implementing the robust counterpart}
\label{sec:experiment}

The preceding section was written, in the first version of this report, as a
theory-only contribution. This section removes that restriction: it implements
the robust placement MILP of \cref{eq:bs-milp} and evaluates it under the same
temporal protocol as \cref{sec:empirical}, on iscf480. Three questions drive
the design. First, is the counterpart \emph{implementable at all} under the
project's own capacity convention --- the $\sigma = 1.10$ slack rule that
calibrates $T_s$? Second, when it is implementable, what does the protection
buy out of sample, and at what price in visits? Third, does \emph{targeted}
budgeted protection beat the obvious nominal alternative --- uniformly
tightening the training ceiling by the same amount of visit price?

\subsection{Setup}
\label{sec:exp-setup}

\paragraph{Solver.}
No commercial solver is available on the experiment machine, so
\cref{eq:bs-milp} is implemented on the open-source HiGHS solver
(\texttt{src/milp\_highs\_robust.py}, via \texttt{highspy}). The visit
variables are declared continuous, which is exact: for fixed integral
$x_{ps}$ the minimal $y_{os} = \max_{p \in P_o} x_{ps}$ is automatically
binary, so the model carries only the $|P|\,|S| = 3{,}840$ assignment binaries
plus the continuous dual block of \cref{prop:bs-counterpart}. Exact
branch-and-bound is nonetheless out of reach at useful time scales: the eight
stations of iscf480 are identical, and the resulting symmetry left HiGHS
unable to improve --- or even find --- an incumbent within minutes. Incumbents
therefore come from a three-stage pipeline: (i) a feasibility-first greedy
(best-fit decreasing on the \emph{robust} station load, swap repair); (ii) a
deterministic pairwise-swap local search on the weighted visit objective, with
the \emph{original} robust constraint --- nominal load plus the sum of the
$\Gamma$ largest deviations, per station --- as a hard accept filter (three
seeded restarts, pass-budgeted for cross-process reproducibility); (iii) a
warm-started HiGHS run used for what an exact solver is still good for here:
\emph{proofs of infeasibility} for budgets that do not fit, and dual bounds.
The optimization objective aggregates the top 1{,}500 weighted distinct
multi-item supports of the training orders; all reported visit counts,
however, are evaluated on the \emph{full real} order sets, and all feasibility
metrics on the fixed $1.10$-rule contract
$T^{\mathrm n} = \lceil 1.10 \sum_p L^{te}_p / (V\,|S|) \rceil$ of
\cref{sec:empirical}. Every solved robust layout was re-verified against the
pre-dualization constraint by direct top-$\Gamma$ arithmetic
(\texttt{robust\_ok} in the result tables): 21 of 21 feasible robust arms
pass, so the dual block never ``cheats.''

\paragraph{Deviations: ex-ante M3 with partial insurance.}
The deviations $\hat L_p$ are calibrated with the only model of
\cref{sec:calibration} available at decision time: M3, the per-product
ddof-$1$ standard deviation of the five volume-rescaled temporal train-block
line counts, capped at $\bar L_p$. (The global-$c$ M2 fit degenerates here:
temporally concentrated SKUs push the $\bar L$-weighted quantile to
$c \approx 70$--$97$, at which point the cap binds for essentially the whole
catalogue and the set collapses to the Soyster box.) Because
\cref{sec:gamma} showed M3 \emph{underestimates} realized drift on iscf480,
and because --- as \cref{sec:exp-impl} shows --- even M3-sized deviations do
not fit the slack contract in full, the experiment adds a partial-insurance
factor $\alpha \in \{0.5, 0.25\}$: the robust constraint is built with
$\alpha \hat L_p$ in place of $\hat L_p$. Formally this is just
\cref{eq:bs-milp} on a shrunken uncertainty set, so
\cref{prop:bs-counterpart} applies verbatim; operationally it asks what
protection is \emph{purchasable} inside the existing capacity margin.

\paragraph{Arms and pairing.}
For each of the five expanding temporal cuts of \cref{sec:empirical}:
the nominal arm $\Gamma = 0$; budgeted arms
$\Gamma \in \{1, 2, 4, 8\} \times \alpha \in \{0.5, 0.25\}$; full-$\sigma$
probes $\Gamma \in \{1, 4, 8, 16\}$ at $\alpha = 1$ (cut~0); two
\emph{uniform-tightening} competitors, trained with the ceiling
$T_s(\beta) = \lceil \beta \sum_p \bar L_p / (V |S|) \rceil$ for
$\beta \in \{1.05, 1.02\}$ and $\Gamma = 0$ --- the same visit price paid as
undirected balance instead of directed protection; and the stored Hexaly
$k{=}1$ layout as an external reference. Every arm of a cut shares the same
supports, budgets, calibration and evaluation, so within-cut comparisons are
paired; the price of robustness $\mathrm{PoR} = (V_{tr}(\cdot) -
V_{tr}(\Gamma{=}0))/V_{tr}(\Gamma{=}0)$ and the test-visit gain
$G = (V_{te}(\Gamma{=}0) - V_{te}(\cdot))/V_{te}(\Gamma{=}0)$ are measured
against the paired nominal arm.

\subsection{Full-\texorpdfstring{$\sigma$}{sigma} protection does not fit the slack contract}
\label{sec:exp-impl}

The $1.10$ rule leaves each station a nominal headroom of
$(1 - 1/1.10)\,T_s$ --- about $4{,}124$ lines per station on cut~0
($T_s = 45{,}365$). The M3 deviations are far larger than that: the top SKU
alone carries $\hat L_p = 13{,}307$ (its block-$\sigma$ exceeds its own
nominal, so the cap binds), and the eight largest deviations sum to
$48{,}621$ --- one and a half times the \emph{total} headroom
$8 \times 4{,}124 \approx 33{,}000$ of the whole cell. Since the heavy
deviators are also heavy products, the nominal constraint prevents parking
them all in one station, and full-$\sigma$ protection has, arithmetically,
nowhere to live. HiGHS makes this formal at moderate budgets:
at $\alpha = 1$ it returns a \textbf{proven infeasibility} certificate for
$\Gamma = 4$, $8$ and $16$ on cut~0, within a minute per arm ($28$--$41$\,s
of total arm time including the constructive pipeline). The
$\Gamma = 1$ case is harder --- its LP relaxation is feasible --- and remains
\textbf{undetermined}: no incumbent and no certificate after a $318$-second
probe and a separate $820$-second attempt (solver logs). The same caution applies at
$\alpha < 1$: the arms marked infeasible in \cref{tab:exp-summary}
($\Gamma \ge 4$ at $\alpha = 0.5$, $\Gamma = 8$ at both $\alpha$, and
$\Gamma = 2$ at $\alpha = 0.5$ on four of five cuts) are
\emph{undetermined-within-budget} --- neither the constructive pipeline nor a
$180$-second HiGHS probe produced a layout --- not proven infeasible. At
half-$\sigma$, by contrast, the top-eight deviation mass is $24{,}310 <
33{,}000$, and indeed $\Gamma = 1$, $\alpha = 0.5$ is constructively feasible
on all five cuts.

The practical reading is blunt: under the project's own capacity convention,
the drift that iscf480 actually exhibits (\cref{sec:gamma}) is too large to be
insured in full. Either the protection target is scaled down ($\alpha$), or
the slack $\sigma$ --- a capacity investment --- must grow. The remainder of
the section quantifies the first option.

\subsection{Results}
\label{sec:exp-results}

\begin{table}[htbp]
\centering\small
\caption{Per-arm means over the temporal cuts on which the arm is feasible
(iscf480; evaluation on real test orders against the fixed $1.10$ contract
$T^{\mathrm n}$). PoR and $G$ are paired against the $\Gamma{=}0$ arm of the
same cut; $G < 0$ means the arm pays visits out of sample too. Arms absent
from the body: at $\alpha=1$, $\Gamma \in \{4,8,16\}$ \emph{proven} infeasible
(HiGHS certificates, cut~0) and $\Gamma = 1$ undetermined; at $\alpha=0.5$,
$\Gamma = 4, 8$ undetermined on 5/5 cuts; at $\alpha = 0.25$, $\Gamma = 8$
undetermined on 5/5 cuts.}
\label{tab:exp-summary}
\begin{tabular}{lrrrrrrr}
\toprule
arm & cuts & PoR\,\% & $G$\,\% & viol.\ ($/8$) & peak $W_s/T^{\mathrm n}$ & worst peak & $X$\,\% \\
\midrule
$\Gamma{=}0$ (nominal)            & 5/5 & 0.00  & 0.00   & 2.6 & 1.77 & 2.11 & 17.4 \\
$\Gamma{=}1$, $\alpha{=}0.5$      & 5/5 & 2.34  & $-$4.44  & 2.6 & \textbf{1.59} & 1.87 & \textbf{13.0} \\
$\Gamma{=}1$, $\alpha{=}0.25$     & 5/5 & 1.44  & $-$2.49  & 2.6 & 1.67 & 1.90 & 15.4 \\
$\Gamma{=}2$, $\alpha{=}0.25$     & 5/5 & 1.41  & $-$1.64  & 2.2 & 1.73 & 1.94 & 15.5 \\
$\Gamma{=}4$, $\alpha{=}0.25$     & 5/5 & 1.77  & $-$2.41  & 2.8 & 1.74 & 2.11 & 15.2 \\
$\Gamma{=}2$, $\alpha{=}0.5$      & 1/5 & 19.13 & $-$23.14 & 3.0 & 1.14 & 1.14 & 3.1 \\
tightened $\beta{=}1.05$          & 5/5 & 0.41  & $-$0.28  & 2.8 & 1.76 & 2.03 & 16.8 \\
tightened $\beta{=}1.02$          & 5/5 & 2.26  & $-$1.73  & 2.8 & 1.75 & 2.01 & 16.9 \\
Hexaly $k{=}1$ (reference)        & 5/5 & $-$3.88 & 6.16 & 2.8 & 2.15 & 2.52 & 23.1 \\
\bottomrule
\end{tabular}
\end{table}

\begin{figure}[htbp]
\centering
\includegraphics[width=\textwidth]{fig_exp_tradeoff}
\caption{Trade-off along the budget $\Gamma$ (means over cuts; open markers =
arms feasible on fewer than five cuts, annotated with their coverage). Left:
price of robustness in train visits. Center: stations above the normalized
ceiling. Right: peak load ratio; the dashed line is contract feasibility.
Partial insurance moves the peak, not the violation count.}
\label{fig:exp-tradeoff}
\end{figure}

\begin{figure}[htbp]
\centering
\includegraphics[width=0.72\textwidth]{fig_exp_frontier}
\caption{Price--severity frontier: mean out-of-sample peak ratio versus PoR.
Budgeted arms (vermillion) move down-right; uniform tightening (blue squares)
moves right \emph{without} moving down; the Hexaly reference (green) shows
that a stronger visit optimizer concentrates load harder. The open marker is
the single-cut $\Gamma{=}2$, $\alpha{=}0.5$ arm.}
\label{fig:exp-frontier}
\end{figure}

\begin{figure}[htbp]
\centering
\includegraphics[width=\textwidth]{fig_exp_stations}
\caption{Per-station load ratios $W_s/T^{\mathrm n}$ on temporal cut~4 (the
worst nominal peak), by arm. The budget compresses the top of the
distribution --- the station a drifting SKU overloads --- rather than
shifting all stations uniformly.}
\label{fig:exp-stations}
\end{figure}

\paragraph{(i) Cheap partial insurance works on severity.}
The efficient point of the sweep is $\Gamma = 1$ at $\alpha = 0.5$: for a
$2.34\%$ visit price it cuts the mean out-of-sample peak from $1.77$ to
$1.59$ and the excess workload fraction $X$ from $17.4\%$ to $13.0\%$
(\cref{tab:exp-summary}). The reduction is \emph{paired sign-consistent}: the
peak ratio falls in five cuts out of five ($1.49\!\to\!1.47$,
$1.48\!\to\!1.34$, $1.81\!\to\!1.61$, $1.95\!\to\!1.87$, $2.11\!\to\!1.64$),
as does $X$. Halving the insured magnitude again ($\alpha = 0.25$) roughly
halves both the price and the effect.

\paragraph{(ii) Equal price, uniform tightening buys nothing.}
The $\beta = 1.02$ competitor pays essentially the same price as the
efficient budgeted arm ($2.26\%$ vs.\ $2.34\%$) and leaves the mean peak at
$1.75$ --- indistinguishable from the nominal $1.77$ --- with per-cut changes
that are sign-\emph{inconsistent} (the peak worsens on two of five cuts).
The mechanism is visible in \cref{fig:exp-stations}: tightening rebalances
\emph{nominal} loads, but the out-of-sample peak is created by large
deviators, and a uniformly tighter ceiling does nothing to separate them ---
the budget constraint does exactly that, because co-locating two large
$\hat L_p$ makes the station's top-$\Gamma$ protection unaffordable. On this
instance family, at equal visit price, \textbf{targeted budgeted protection
dominates uniform tightening on peak severity}.

\paragraph{(iii) The feasibility wall is sharp.}
One step further, $\Gamma = 2$ at $\alpha = 0.5$, is constructible on a
single cut --- and there it nearly restores contract feasibility (peak
$1.14$, $X = 3.1\%$) at a $19.1\%$ visit price. Insurance beyond the slack
headroom is not gradually more expensive; it is discontinuously so, which is
the operational face of the arithmetic in \cref{sec:exp-impl}.

\paragraph{(iv) Severity, not incidence.}
Violation \emph{counts} barely move across implementable arms
($2.2$--$2.8$ of $8$; the count can even rise while the peak falls, as on
cut~3). This is exactly what \cref{sec:gamma} predicts: M3 underestimates the
realized drift, so ex-ante-calibrated protection trims how far stations
overshoot, not whether some station somewhere overshoots. Restoring
feasibility outright would need the drift-scale magnitudes
($c_1(0.9) \approx 1.2$--$1.6$) that no train-internal model provides.

\paragraph{(v) The visit--fragility trade-off, priced.}
The Hexaly reference beats our incumbents by $3.88\%$ on train visits (and
$6.16\%$ on test) --- and carries the worst workload profile of the table
(peak $2.15$, $X = 23.1\%$). Better visit optimization concentrates
correlated, drifting SKUs; the budget constraint is the first dial in this
project that prices that concentration explicitly instead of discovering it
out of sample.

\subsection{Solver honesty}
\label{sec:exp-honesty}

Three caveats bound the claims. First, incumbents are metaheuristic: on the
truncated support objective of cut~0 the HiGHS dual bound is $30{,}312$
against an incumbent of $43{,}158$ ($\approx 30\%$ gap, largely inherited
from the weak LP relaxation of the visit objective), so PoR magnitudes carry
solver noise --- the paired, sign-consistent within-cut deltas are the
reliable signal, not third digits. Second, HiGHS wall-times overran its
\texttt{time\_limit} option on this model (up to $820$\,s against a
$480$\,s limit in the solver logs; persisted arm times reach $465$\,s
against $180$\,s probes), which is why the sweep runs on the deterministic
local-search pipeline and reserves HiGHS for certificates and bounds. Third,
$G < 0$ throughout: every robust arm pays its visit price out of sample as
well ($-1.6\%$ to $-4.4\%$ at the implementable budgets). Nothing here claims
a visits benefit --- the insured quantity is workload-feasibility risk, in
line with the falsified-H1 lesson that opened this report.


\section{Discussion and recommendations}
\label{sec:discussion}

\paragraph{When is the robust workload constraint worth it?}
The evidence separates three regimes.
(i) \emph{Stationary demand, matched volume} (iscf480 CV): the nominal constraint is
adequate --- normalized violations exist only because placement sits at the exact slack
edge, and they are immaterial (peak ratio $\le 1.016$, $X \le 0.40\%$). Robustness would
buy little here.
(ii) \emph{Temporal drift} (all temporal families, and the planted demand shift): the
nominal constraint fails materially and systematically --- 1--4 of 8 stations, peak ratios
up to 2.52 on iscf480 and 2.63 on the full-catalogue $k$ grid (3.45 in the sparse-train
variants), 7--37\% of workload above the fair-share ceiling --- and the failure is invisible
to the as-run check whenever $B < 1/\sigma$. This is the regime the counterpart of
\cref{sec:theory} targets, and the calibration says protection is cheap in budget terms:
the realized overloads are covered at median $\Gamma^{\mathrm{cover}}$ of 8/60 and 19/541
products under the best-fitting magnitude model --- strikingly close to the
$\Gamma \sim \sqrt{n}$ rule of thumb of \cref{rem:bs-prob}
($\sqrt{60} \approx 7.7$, $\sqrt{541} \approx 23.3$) --- and deviations are
head-concentrated (Gini 0.63/0.82), the Bertsimas--Sim sweet spot.
(iii) \emph{Volume-artifact regimes} ($B \gg 1$ or $B \ll 1/\sigma$): fix the protocol
first; no constraint reformulation can rescue a check that compares incommensurate volumes.

\paragraph{What $\Gamma$ and what magnitudes the evidence suggests.}
A per-station budget of order $\sqrt{\zeta_s}$ (reported as the fraction
$\Gamma_s / \zeta_s \approx 13\%$ and $4\%$ on the two families) would have covered the
median realized overload. Because $\Gamma^{\mathrm{bind}} = 1$ at 58--68\% of stations,
even $\Gamma = 1$ already reshapes the training solution --- the practical risk is
over-conservatism, not vacuity, which argues for starting at small budgets and the
fractional-budget refinement rather than box protection. The harder open problem is the
\emph{magnitude} side: the required $c_1(0.9) \approx 1.2$--$1.6$ (growing with horizon)
is far beyond what train-internal variability (M3) predicts on iscf480, so an ex-ante
deviation model needs longer demand history or an explicit drift model, not just
within-window dispersion; on the full-catalogue instance M3 is a usable starting point
(38--80\% coverage).

\paragraph{What the implementation shows.}
The computational study of \cref{sec:experiment} turns the counterpart from a proposal
into three practical statements. (a) \emph{Partial ex-ante insurance is cheap and
effective on severity}: $\Gamma = 1$ at half the M3 magnitude cuts the out-of-sample peak
by ten percent of its value and the excess workload by a quarter, for a
$2.3\%$ visit price, sign-consistently across all five cuts. (b) \emph{Uniform
tightening is not a substitute}: at the same visit price it leaves the peak where it was,
because it rebalances nominal loads without separating the large deviators that create
the peak --- the budget constraint does that separation by construction. (c)
\emph{Full-magnitude protection is a capacity decision, not a modeling one}: the top
M3 deviations alone exceed the whole cell's slack headroom ($48.6$k vs.\ $33.0$k lines
on the first cut), HiGHS certifies infeasibility from $\Gamma = 4$ up, and the one
constructible step beyond the headroom ($\Gamma = 2$, $\alpha = 0.5$, one cut) costs
$19\%$ of train visits. Insuring the drift this data actually exhibits requires growing
the slack $\sigma$ --- a capacity investment that \cref{sec:exp-impl} prices.

\paragraph{Threats to validity.}
(i) The five temporal cuts share train mass (expanding windows): their spread is a
dispersion proxy, not a confidence interval, and per-fold results are correlated.
(ii) The volume-normalized ceiling is a \emph{counterfactual protocol choice}
($T^{\mathrm n}$ re-solves nothing; it asks how the trained layout would fare if the
ceiling had been calibrated on the test volume). Other normalizations are defensible; the
$B$ diagnosis of the as-run check does not depend on this choice.
(iii) Layouts are Hexaly (metaheuristic) placements: balancing exactly to the slack edge
is a solver behavior, and better or different placements could shift the borderline CV
findings (though hardly the 1.7--2.5$\times$ temporal overloads).
(iv) Instance breadth: all synthetic families derive from one warehouse's catalogue, and
the industrial evidence is a single BERNER cell with a single temporal cut and no persisted
layouts (its bindingness verdict rests on assignment bounds).
(v) The $\Gamma$ calibration is retrospective by design; only M3 is available ex ante, and
it underestimates drift on iscf480 by an order of magnitude.
(vi) The experiment's incumbents are metaheuristic, with a $\approx 30\%$ HiGHS dual-bound
gap on the truncated support objective; paired within-cut deltas are reliable, absolute
PoR magnitudes less so.
(vii) The partial-insurance grid is coarse ($\alpha \in \{0.5, 0.25\}$), and the
infeasible robust arms at $\alpha < 1$ are undetermined-within-budget rather than proven
infeasible; only the full-magnitude certificates at $\Gamma \in \{4, 8, 16\}$ are proofs.

\appendix

\section{Full per-fold tables}
\label{app:tables}

\begin{table}[htbp]
\centering
\caption{iscf480, all stored (split, fold, $k$) cells. \texttt{asrun} = stored
\texttt{wl\_broken} (reproduced exactly); ratio$^{\mathrm n}$ = peak
$W^{\mathrm{te}}_s / T^{\mathrm n}_s$; imb.\ = imbalance. Source:
\texttt{results/iscf480/per\_fold.csv}.}
\label{tab:app480}
\small
\begin{tabular}{@{}llccccccc@{}}
\toprule
Split & Fold & $k$ & as-run & norm & ratio$^{\mathrm n}$ & imb. & $B$ & $X$ (\%) \\
\midrule
CV & r0f0 & 1 & 0 & \textbf{2} & 1.016 & 1.117 & 0.228 & 0.24 \\
CV & r0f0 & 2 & 0 & \textbf{2} & 1.014 & 1.116 & 0.228 & 0.25 \\
CV & r0f1 & 1 & 0 & 0 & 0.999 & 1.098 & 0.227 & 0.00 \\
CV & r0f1 & 2 & 0 & 0 & 0.994 & 1.094 & 0.227 & 0.00 \\
CV & r0f2 & 1 & 0 & \textbf{2} & 1.005 & 1.105 & 0.227 & 0.13 \\
CV & r0f2 & 2 & 0 & \textbf{1} & 1.006 & 1.106 & 0.227 & 0.08 \\
CV & r0f3 & 1 & 0 & \textbf{3} & 1.012 & 1.114 & 0.227 & 0.40 \\
CV & r0f3 & 2 & 0 & \textbf{1} & 1.004 & 1.105 & 0.227 & 0.06 \\
CV & r0f4 & 1 & 0 & \textbf{1} & 1.000 & 1.100 & 0.227 & 0.00 \\
CV & r0f4 & 2 & 0 & 0 & 0.999 & 1.099 & 0.227 & 0.00 \\
\midrule
temp. & r0f0 & 1 & \textbf{4} & \textbf{3} & 1.714 & 1.886 & 1.005 & 11.02 \\
temp. & r0f0 & 2 & \textbf{2} & \textbf{2} & 1.541 & 1.696 & 1.005 & 13.37 \\
temp. & r0f1 & 1 & \textbf{2} & \textbf{3} & 1.906 & 2.097 & 0.693 & 19.81 \\
temp. & r0f1 & 2 & \textbf{1} & \textbf{4} & 1.710 & 1.881 & 0.693 & 14.75 \\
temp. & r0f2 & 1 & \textbf{1} & \textbf{3} & 2.255 & 2.480 & 0.457 & 26.61 \\
temp. & r0f2 & 2 & 0 & \textbf{3} & 1.981 & 2.179 & 0.457 & 21.19 \\
temp. & r0f3 & 1 & 0 & \textbf{2} & 2.519 & 2.771 & 0.281 & 27.23 \\
temp. & r0f3 & 2 & 0 & \textbf{3} & 2.189 & 2.408 & 0.281 & 28.37 \\
temp. & r0f4 & 1 & 0 & \textbf{3} & 2.376 & 2.614 & 0.125 & 30.72 \\
temp. & r0f4 & 2 & 0 & \textbf{3} & 1.882 & 2.070 & 0.125 & 30.25 \\
\bottomrule
\end{tabular}
\end{table}

\begin{table}[htbp]
\centering
\caption{iscf10kt temporal variant, full $k$ grid. Source:
\texttt{results/iscf10kt/per\_fold.csv}.}
\label{tab:app10kt}
\small
\begin{tabular}{@{}lccccccc@{}}
\toprule
Cut & $k$ & as-run & norm & ratio$^{\mathrm n}$ & imb. & $B$ & $X$ (\%) \\
\midrule
0 & 1 & \textbf{2} & \textbf{2} & 1.50 & 1.65 & 0.79 & 10.6 \\
0 & 2 & \textbf{2} & \textbf{2} & 1.47 & 1.61 & 0.79 & 11.8 \\
0 & 3 & \textbf{1} & \textbf{2} & 1.56 & 1.71 & 0.79 & 8.6 \\
0 & 4 & \textbf{2} & \textbf{2} & 1.32 & 1.46 & 0.79 & 8.3 \\
0 & 6 & \textbf{1} & \textbf{1} & 1.49 & 1.64 & 0.79 & 6.8 \\
1 & 1 & \textbf{1} & \textbf{2} & 1.69 & 1.86 & 0.55 & 15.6 \\
1 & 2 & 0 & \textbf{2} & 1.53 & 1.69 & 0.55 & 11.8 \\
1 & 3 & 0 & \textbf{1} & 1.56 & 1.72 & 0.55 & 7.7 \\
1 & 4 & \textbf{1} & \textbf{2} & 1.65 & 1.82 & 0.55 & 9.9 \\
1 & 6 & 0 & \textbf{3} & 1.44 & 1.58 & 0.55 & 10.5 \\
2 & 1 & 0 & \textbf{1} & 2.15 & 2.36 & 0.37 & 15.8 \\
2 & 2 & 0 & \textbf{2} & 1.66 & 1.83 & 0.37 & 14.2 \\
2 & 3 & 0 & \textbf{2} & 2.06 & 2.27 & 0.37 & 19.2 \\
2 & 4 & 0 & \textbf{2} & 1.90 & 2.09 & 0.37 & 13.0 \\
2 & 6 & 0 & \textbf{3} & 1.75 & 1.92 & 0.37 & 13.0 \\
3 & 1 & 0 & \textbf{3} & 1.93 & 2.12 & 0.22 & 24.6 \\
3 & 2 & 0 & \textbf{1} & 2.36 & 2.60 & 0.22 & 18.8 \\
3 & 3 & 0 & \textbf{2} & 2.16 & 2.38 & 0.22 & 24.5 \\
3 & 4 & 0 & \textbf{3} & 1.78 & 1.96 & 0.22 & 22.6 \\
3 & 6 & 0 & \textbf{3} & 1.68 & 1.85 & 0.22 & 20.0 \\
4 & 1 & 0 & \textbf{3} & 2.35 & 2.59 & 0.11 & 25.4 \\
4 & 2 & 0 & \textbf{2} & 2.63 & 2.90 & 0.11 & 36.8 \\
4 & 3 & 0 & \textbf{2} & 2.26 & 2.49 & 0.11 & 28.1 \\
4 & 4 & 0 & \textbf{3} & 2.55 & 2.81 & 0.11 & 28.2 \\
4 & 6 & 0 & \textbf{3} & 2.03 & 2.23 & 0.11 & 22.2 \\
\bottomrule
\end{tabular}
\end{table}

\begin{table}[htbp]
\centering
\caption{iscf10kt sparse-train variants, $k{=}1$ rows (full grid, $k \in \{2,3,4,6\}$
included, on disk). Source: \texttt{results/iscf10kt/per\_fold.csv}.}
\label{tab:appsparse}
\small
\begin{tabular}{@{}llcccccc@{}}
\toprule
Variant & Fold & as-run & norm & ratio$^{\mathrm n}$ & imb. & $B$ & $X$ (\%) \\
\midrule
t500  & 0 & \textbf{8} & \textbf{1} & 1.78 & 1.96 & 8.18 & 10.8 \\
t500  & 1 & \textbf{8} & \textbf{1} & 1.90 & 2.09 & 6.61 & 12.4 \\
t500  & 2 & \textbf{8} & \textbf{1} & 2.30 & 2.53 & 4.91 & 17.8 \\
t500  & 3 & \textbf{8} & \textbf{3} & 2.14 & 2.36 & 3.50 & 19.8 \\
t500  & 4 & \textbf{4} & \textbf{2} & 3.45 & 3.80 & 2.07 & 35.5 \\
t1000 & 0 & \textbf{8} & \textbf{2} & 1.56 & 1.72 & 4.03 & 11.1 \\
t1000 & 1 & \textbf{8} & \textbf{1} & 1.84 & 2.02 & 3.28 & 11.5 \\
t1000 & 2 & \textbf{8} & \textbf{2} & 2.09 & 2.30 & 2.53 & 15.2 \\
t1000 & 3 & \textbf{6} & \textbf{3} & 1.71 & 1.89 & 1.77 & 18.4 \\
t1000 & 4 & \textbf{2} & \textbf{2} & 3.44 & 3.79 & 1.06 & 35.4 \\
t2000 & 0 & \textbf{7} & \textbf{2} & 1.54 & 1.70 & 1.97 & 9.2 \\
t2000 & 1 & \textbf{7} & \textbf{3} & 1.49 & 1.63 & 1.68 & 13.2 \\
t2000 & 2 & \textbf{4} & \textbf{3} & 1.68 & 1.85 & 1.27 & 15.6 \\
t2000 & 3 & \textbf{3} & \textbf{3} & 2.08 & 2.29 & 0.90 & 22.5 \\
t2000 & 4 & \textbf{1} & \textbf{2} & 2.69 & 2.96 & 0.52 & 31.0 \\
\bottomrule
\end{tabular}
\end{table}

\begin{table}[htbp]
\centering
\caption{iscfco base arms (unshifted $\rho = 0$ evaluation). $k{=}1$ metrics are
identical across the mm/ms variants by construction. Per-fold $B$ = 0.858, 0.573, 0.374,
0.221, 0.099 for cuts 0--4. Source: \texttt{results/secondary/per\_fold.csv}.}
\label{tab:appco}
\small
\begin{tabular}{@{}llcccccc@{}}
\toprule
Variant & Cut & $k$ & as-run & norm & ratio$^{\mathrm n}$ & imb. & $X$ (\%) \\
\midrule
mm & 0 & 1  & 0 & \textbf{3} & 1.059 & 1.165 & 1.46 \\
mm & 0 & 10 & 0 & \textbf{1} & 1.028 & 1.131 & 0.39 \\
mm & 1 & 1  & 0 & \textbf{2} & 1.055 & 1.161 & 1.38 \\
mm & 1 & 10 & 0 & \textbf{2} & 1.055 & 1.160 & 1.20 \\
mm & 2 & 1  & 0 & \textbf{2} & 1.045 & 1.150 & 1.20 \\
mm & 2 & 10 & 0 & \textbf{2} & 1.060 & 1.167 & 1.17 \\
mm & 3 & 1  & 0 & \textbf{2} & 1.089 & 1.198 & 1.27 \\
mm & 3 & 10 & 0 & \textbf{2} & 1.064 & 1.171 & 1.24 \\
mm & 4 & 1  & 0 & \textbf{3} & 1.057 & 1.164 & 1.96 \\
mm & 4 & 10 & 0 & \textbf{3} & 1.083 & 1.193 & 1.52 \\
ms & 0 & 1  & 0 & \textbf{3} & 1.059 & 1.165 & 1.46 \\
ms & 0 & 10 & \textbf{1} & \textbf{1} & 1.061 & 1.167 & 0.83 \\
ms & 1 & 1  & 0 & \textbf{2} & 1.055 & 1.161 & 1.38 \\
ms & 1 & 10 & 0 & \textbf{2} & 1.064 & 1.171 & 1.15 \\
ms & 2 & 1  & 0 & \textbf{2} & 1.045 & 1.150 & 1.20 \\
ms & 2 & 10 & 0 & \textbf{1} & 1.014 & 1.116 & 0.20 \\
ms & 3 & 1  & 0 & \textbf{2} & 1.089 & 1.198 & 1.27 \\
ms & 3 & 10 & 0 & \textbf{3} & 1.045 & 1.149 & 0.91 \\
ms & 4 & 1  & 0 & \textbf{3} & 1.057 & 1.164 & 1.96 \\
ms & 4 & 10 & 0 & \textbf{2} & 1.089 & 1.199 & 1.50 \\
\bottomrule
\end{tabular}
\end{table}

\begin{table}[htbp]
\centering
\caption{Planted-affinity informative arms (reps r0, r2; rep r1 is degenerate as-run with
empty layouts and is omitted). The as-run 7--8/8 counts are volume artifact ($B > 1$).
Source: \texttt{results/secondary/per\_fold.csv}.}
\label{tab:appaff}
\small
\begin{tabular}{@{}lccccccc@{}}
\toprule
Arm & $k$ & as-run & norm & ratio$^{\mathrm n}$ & imb. & $B$ & $X$ (\%) \\
\midrule
combine\ $\alpha{=}0.0$\ r0 & 1  & \textbf{8} & \textbf{2} & 1.071 & 1.180 & 1.81 & 1.93 \\
combine\ $\alpha{=}0.0$\ r0 & 10 & \textbf{8} & \textbf{1} & 1.065 & 1.173 & 1.81 & 0.89 \\
combine\ $\alpha{=}0.5$\ r0 & 1  & \textbf{8} & \textbf{2} & 1.049 & 1.155 & 3.18 & 1.05 \\
combine\ $\alpha{=}0.5$\ r0 & 10 & \textbf{8} & \textbf{2} & 1.052 & 1.158 & 3.18 & 0.88 \\
combine\ $\alpha{=}1.0$\ r0 & 1  & \textbf{8} & \textbf{2} & 1.063 & 1.170 & 4.40 & 1.25 \\
combine\ $\alpha{=}1.0$\ r0 & 10 & \textbf{8} & \textbf{3} & 1.055 & 1.161 & 4.40 & 1.83 \\
reassign\ $\alpha{=}0.0$\ r0 & 1  & \textbf{8} & \textbf{2} & 1.092 & 1.202 & 1.82 & 2.42 \\
reassign\ $\alpha{=}0.0$\ r0 & 10 & \textbf{8} & \textbf{3} & 1.093 & 1.203 & 1.82 & 1.89 \\
reassign\ $\alpha{=}0.5$\ r0 & 1  & \textbf{8} & \textbf{2} & 1.170 & 1.288 & 1.93 & 3.43 \\
reassign\ $\alpha{=}0.5$\ r0 & 10 & \textbf{8} & \textbf{2} & 1.180 & 1.298 & 1.93 & 4.51 \\
reassign\ $\alpha{=}1.0$\ r0 & 1  & \textbf{8} & \textbf{1} & 1.165 & 1.282 & 1.79 & 2.27 \\
reassign\ $\alpha{=}1.0$\ r0 & 10 & \textbf{8} & \textbf{2} & 1.196 & 1.317 & 1.79 & 3.68 \\
combine\ $\alpha{=}0.0$\ r2 & 1  & \textbf{8} & \textbf{3} & 1.130 & 1.243 & 1.85 & 2.62 \\
combine\ $\alpha{=}0.0$\ r2 & 10 & \textbf{7} & \textbf{2} & 1.129 & 1.242 & 1.85 & 2.03 \\
combine\ $\alpha{=}0.5$\ r2 & 1  & \textbf{8} & \textbf{2} & 1.039 & 1.144 & 3.26 & 0.79 \\
combine\ $\alpha{=}0.5$\ r2 & 10 & \textbf{8} & \textbf{3} & 1.071 & 1.179 & 3.26 & 1.68 \\
combine\ $\alpha{=}1.0$\ r2 & 1  & \textbf{8} & \textbf{2} & 1.064 & 1.171 & 4.48 & 1.12 \\
combine\ $\alpha{=}1.0$\ r2 & 10 & \textbf{8} & \textbf{3} & 1.035 & 1.139 & 4.48 & 1.13 \\
reassign\ $\alpha{=}0.0$\ r2 & 1  & \textbf{8} & \textbf{2} & 1.128 & 1.243 & 1.84 & 2.04 \\
reassign\ $\alpha{=}0.0$\ r2 & 10 & \textbf{8} & \textbf{2} & 1.124 & 1.237 & 1.84 & 2.00 \\
reassign\ $\alpha{=}0.5$\ r2 & 1  & \textbf{8} & \textbf{2} & 1.196 & 1.317 & 1.88 & 3.84 \\
reassign\ $\alpha{=}0.5$\ r2 & 10 & \textbf{7} & \textbf{4} & 1.244 & 1.370 & 1.88 & 7.40 \\
reassign\ $\alpha{=}1.0$\ r2 & 1  & \textbf{8} & \textbf{3} & 1.150 & 1.265 & 1.90 & 2.78 \\
reassign\ $\alpha{=}1.0$\ r2 & 10 & \textbf{8} & \textbf{3} & 1.158 & 1.274 & 1.90 & 5.54 \\
\bottomrule
\end{tabular}
\end{table}

\begin{table}[htbp]
\centering
\caption{Computational study (\cref{sec:experiment}): all feasible (cut, arm)
cells on iscf480 temporal. PoR and $G$ are paired against the cut's
$\Gamma{=}0$ arm; viol.\ = stations above $T^{\mathrm n}$ (of 8); peak =
$\max_s W^{\mathrm{te}}_s / T^{\mathrm n}$. Arms not shown are
undetermined-within-budget ($\Gamma{=}2$, $\alpha{=}0.5$ on cuts 1--4;
$\Gamma{=}4, 8$ at $\alpha{=}0.5$ and $\Gamma{=}8$ at $\alpha{=}0.25$ on all
cuts; $\Gamma{=}1$ at $\alpha{=}1$ on cut 0) or proven infeasible
($\alpha{=}1$, $\Gamma \in \{4,8,16\}$, cut 0).
Source: \texttt{results/experiment/results.csv}.}
\label{tab:app-exp}
\small
\begin{tabular}{@{}llrrrrr@{}}
\toprule
Cut & Arm & PoR (\%) & $G$ (\%) & viol. & peak & $X$ (\%) \\
\midrule
0 & $\Gamma{=}0$ & 0.00 & 0.00 & 2 & 1.489 & 10.6 \\
0 & $\Gamma{=}1$, $\alpha{=}0.5$ & 2.11 & $-$3.49 & 1 & 1.468 & 6.4 \\
0 & $\Gamma{=}1$, $\alpha{=}0.25$ & 1.70 & $-$3.49 & 2 & 1.391 & 8.6 \\
0 & $\Gamma{=}2$, $\alpha{=}0.5$ & 19.13 & $-$23.14 & 3 & 1.139 & 3.1 \\
0 & $\Gamma{=}2$, $\alpha{=}0.25$ & 1.21 & $-$2.47 & 3 & 1.408 & 8.7 \\
0 & $\Gamma{=}4$, $\alpha{=}0.25$ & 2.02 & $-$1.89 & 3 & 1.429 & 9.4 \\
0 & tight $\beta{=}1.05$ & 0.64 & $-$0.60 & 2 & 1.519 & 9.2 \\
0 & tight $\beta{=}1.02$ & 3.10 & $-$3.72 & 2 & 1.498 & 8.3 \\
0 & Hexaly $k{=}1$ & $-$3.45 & 3.41 & 3 & 1.714 & 11.0 \\
\addlinespace
1 & $\Gamma{=}0$ & 0.00 & 0.00 & 3 & 1.482 & 14.0 \\
1 & $\Gamma{=}1$, $\alpha{=}0.5$ & 2.39 & $-$4.06 & 3 & 1.342 & 10.7 \\
1 & $\Gamma{=}1$, $\alpha{=}0.25$ & 1.39 & $-$2.53 & 3 & 1.517 & 10.9 \\
1 & $\Gamma{=}2$, $\alpha{=}0.25$ & 1.66 & $-$1.44 & 2 & 1.710 & 11.5 \\
1 & $\Gamma{=}4$, $\alpha{=}0.25$ & 0.14 & $-$1.78 & 3 & 1.540 & 9.4 \\
1 & tight $\beta{=}1.05$ & 0.32 & 0.25 & 2 & 1.533 & 12.4 \\
1 & tight $\beta{=}1.02$ & 0.91 & $-$1.87 & 2 & 1.427 & 11.0 \\
1 & Hexaly $k{=}1$ & $-$5.05 & 7.09 & 3 & 1.906 & 19.8 \\
\addlinespace
2 & $\Gamma{=}0$ & 0.00 & 0.00 & 3 & 1.814 & 17.4 \\
2 & $\Gamma{=}1$, $\alpha{=}0.5$ & 1.63 & $-$3.02 & 3 & 1.613 & 15.4 \\
2 & $\Gamma{=}1$, $\alpha{=}0.25$ & 0.47 & 0.29 & 3 & 1.793 & 17.3 \\
2 & $\Gamma{=}2$, $\alpha{=}0.25$ & 1.53 & $-$1.91 & 2 & 1.724 & 15.8 \\
2 & $\Gamma{=}4$, $\alpha{=}0.25$ & 1.32 & $-$1.31 & 3 & 1.876 & 15.4 \\
2 & tight $\beta{=}1.05$ & 0.35 & $-$1.05 & 3 & 1.903 & 15.4 \\
2 & tight $\beta{=}1.02$ & 2.28 & $-$0.40 & 3 & 1.935 & 17.2 \\
2 & Hexaly $k{=}1$ & $-$4.18 & 7.22 & 3 & 2.255 & 26.6 \\
\addlinespace
3 & $\Gamma{=}0$ & 0.00 & 0.00 & 2 & 1.948 & 21.4 \\
3 & $\Gamma{=}1$, $\alpha{=}0.5$ & 2.17 & $-$2.69 & 4 & 1.867 & 18.5 \\
3 & $\Gamma{=}1$, $\alpha{=}0.25$ & 1.82 & $-$4.35 & 2 & 1.899 & 19.6 \\
3 & $\Gamma{=}2$, $\alpha{=}0.25$ & 0.90 & $-$0.19 & 2 & 1.868 & 22.1 \\
3 & $\Gamma{=}4$, $\alpha{=}0.25$ & 2.73 & $-$0.21 & 2 & 2.110 & 23.0 \\
3 & tight $\beta{=}1.05$ & $-$0.42 & 2.69 & 3 & 2.034 & 25.2 \\
3 & tight $\beta{=}1.02$ & 1.88 & 1.87 & 3 & 1.901 & 25.2 \\
3 & Hexaly $k{=}1$ & $-$3.30 & 6.27 & 2 & 2.519 & 27.2 \\
\addlinespace
4 & $\Gamma{=}0$ & 0.00 & 0.00 & 3 & 2.112 & 23.5 \\
4 & $\Gamma{=}1$, $\alpha{=}0.5$ & 3.43 & $-$8.94 & 2 & 1.643 & 14.0 \\
4 & $\Gamma{=}1$, $\alpha{=}0.25$ & 1.83 & $-$2.35 & 3 & 1.746 & 20.5 \\
4 & $\Gamma{=}2$, $\alpha{=}0.25$ & 1.76 & $-$2.19 & 2 & 1.940 & 19.6 \\
4 & $\Gamma{=}4$, $\alpha{=}0.25$ & 2.62 & $-$6.87 & 3 & 1.770 & 18.5 \\
4 & tight $\beta{=}1.05$ & 1.17 & $-$2.68 & 4 & 1.818 & 22.0 \\
4 & tight $\beta{=}1.02$ & 3.11 & $-$4.52 & 4 & 2.014 & 23.1 \\
4 & Hexaly $k{=}1$ & $-$3.41 & 6.81 & 3 & 2.376 & 30.7 \\
\bottomrule
\end{tabular}
\end{table}

\begin{thebibliography}{9}
\bibitem{bertsimas2004}
D.~Bertsimas and M.~Sim.
The price of robustness.
\emph{Operations Research}, 52(1):35--53, 2004.

\bibitem{bertsimas2003robust}
D.~Bertsimas and M.~Sim.
Robust discrete optimization and network flows.
\emph{Mathematical Programming}, 98(1--3):49--71, 2003.

\bibitem{soyster1973convex}
A.~L. Soyster.
Convex programming with set-inclusive constraints and applications to inexact linear
programming.
\emph{Operations Research}, 21(5):1154--1157, 1973.

\bibitem{bental2009robust}
A.~Ben-Tal, L.~El~Ghaoui, and A.~Nemirovski.
\emph{Robust Optimization}.
Princeton University Press, Princeton, NJ, 2009.
\end{thebibliography}

\end{document}
